\chapter{Chapter04 — Switching}
%%% generated by the blueprint pipeline from FABL.Chapter04.Switching
\section{Overview}

Declarations of the Lean module \texttt{FABL.Chapter04.Switching}.

\section{Declarations}

\begin{definition}[Binary function under a fixed restriction]
\lean{FABL.restrictedBinaryFunction}
\statementsource{lmn-constant-depth-fourier}{pdf-63759bd909d5-p005-b001}
\statementsource{arora-barak-ch13-circuit-lowerbounds}{pdf-52929dbaa962-p002-b003}
\statementsource{Analysis-of-Boolean-Functions-by-Ryan-ODonnell}{
  pdf-ffbe9c7b3306-p075-b002,
  pdf-ffbe9c7b3306-p098-b004
}
\label{FABL.restrictedBinaryFunction}
\leanok
\uses{FABL.BooleanFunction, FABL.FixedSignCube, FABL.Sign, FABL.binaryCubeSignEquiv, FABL.extendedSignRestriction}
Let $n$ be a nonnegative integer, let $f$ be a function from the sign cube
$\{−1,1\}^n$ to $\{−1,1\}$, let $J\subseteq\{1,\ldots,n\}$, and fix signs on the
coordinates outside $J$. The associated function on the binary cube $\mathbb F_2^n$
sends each binary input to its sign-cube encoding, extends the fixed signs on the
coordinates outside $J$, and evaluates $f$ on the resulting sign vector.
\end{definition}

\begin{definition}[Restricted decision tree depth]
\lean{FABL.restrictedDecisionTreeDepth}
\label{FABL.restrictedDecisionTreeDepth}
\leanok
\uses{FABL.BooleanFunction, FABL.FixedSignCube, FABL.F₂DecisionTree.decisionTreeDepth, FABL.restrictedBinaryFunction}
Let $f$ be a Boolean function on $n$ coordinates, let $J$ be a set of coordinates, and
fix a sign assignment on the coordinates outside $J$. Define the restricted decision
tree depth to be the decision tree depth of the binary function obtained by restricting
$f$ according to these choices.
\end{definition}

\begin{definition}[Switching failure indicator]
\lean{FABL.switchingFailureIndicator}
\statementsource{lmn-constant-depth-fourier}{pdf-63759bd909d5-p005-b004}
\statementsource{arora-barak-ch13-circuit-lowerbounds}{pdf-52929dbaa962-p002-b003}
\label{FABL.switchingFailureIndicator}
\leanok
\uses{FABL.BooleanFunction, FABL.FixedSignCube, FABL.restrictedDecisionTreeDepth}
Let $n$ be a nonnegative integer, let $f$ be a Boolean function on $n$ sign coordinates,
let $k$ be a nonnegative integer, let $J\subseteq\{1,\ldots,n\}$, and fix signs on
the coordinates outside $J$. The switching failure indicator is $1$ if the restricted
decision tree for $f$ has depth at least $k$, and is $0$ otherwise.
\end{definition}

\begin{definition}[Switching failure probability]
\lean{FABL.switchingFailureProbability}
\statementsource{lmn-constant-depth-fourier}{pdf-63759bd909d5-p005-b004}
\statementsource{arora-barak-ch13-circuit-lowerbounds}{pdf-52929dbaa962-p002-b003}
\label{FABL.switchingFailureProbability}
\leanok
\uses{FABL.BooleanFunction, FABL.expectRandomRestriction, FABL.switchingFailureIndicator}
Let $n$ be a nonnegative integer, let $f$ be a Boolean function on $n$ sign coordinates,
let $\delta\in\mathbb R$, and let $k$ be a nonnegative integer. The switching failure
probability is the expected value, under the random restriction with parameter
$\delta$, of the switching failure indicator at threshold $k$.
\end{definition}

\begin{theorem}[Nonnegativity of the switching failure indicator]
\lean{FABL.switchingFailureIndicator_nonneg}
\label{FABL.switchingFailureIndicator_nonneg}
\leanok
\uses{FABL.BooleanFunction, FABL.FixedSignCube, FABL.switchingFailureIndicator}
\difficulty{4}
Let $n,k\in\mathbb N$, let $f:\{−1,1\}^n\to\{−1,1\}$, let $J\subseteq\{1,\ldots,n\}$,
and fix signs on the coordinates outside $J$. Form the restricted function by varying
the signs on $J$ while keeping those outside signs fixed, and let its decision-tree
depth be the least depth of a binary decision tree computing the corresponding function
on $\{0,1\}^n$. The indicator that this depth is at least $k$ is nonnegative.
\end{theorem}

\begin{theorem}[Upper bound for the switching failure indicator]
\lean{FABL.switchingFailureIndicator_le_one}
\label{FABL.switchingFailureIndicator_le_one}
\leanok
\uses{FABL.BooleanFunction, FABL.FixedSignCube, FABL.switchingFailureIndicator}
\difficulty{4}
Let $n,k\in\mathbb N$, let $f:\{−1,1\}^n\to\{−1,1\}$, let $J\subseteq\{1,\ldots,n\}$,
and fix signs on the coordinates outside $J$. Form the restricted function by varying
the signs on $J$ while keeping those outside signs fixed, and let its decision-tree
depth be the least depth of a binary decision tree computing the corresponding function
on $\{0,1\}^n$. The indicator that this depth is at least $k$ is at most $1$.
\end{theorem}

\begin{theorem}[Switching failure indicator at zero threshold]
\lean{FABL.switchingFailureIndicator_zero}
\label{FABL.switchingFailureIndicator_zero}
\leanok
\uses{FABL.BooleanFunction, FABL.FixedSignCube, FABL.switchingFailureIndicator}
\difficulty{1}
Let $n\in\mathbb N$, let $f:\{−1,1\}^n\to\{−1,1\}$, let $J\subseteq\{1,\ldots,n\}$, and
fix signs on the coordinates outside $J$. The indicator that the decision-tree depth of
the restriction of $f$ to those fixed signs is at least zero equals one.
\end{theorem}

\begin{theorem}[Switching failure at threshold zero]
\lean{FABL.switchingFailureProbability_zero}
\label{FABL.switchingFailureProbability_zero}
\leanok
\uses{FABL.BooleanFunction, FABL.expectRandomRestriction, FABL.sum_deltaRandomSubsetWeight, FABL.switchingFailureIndicator, FABL.switchingFailureProbability}
\difficulty{1}
Let $f:\{-1,1\}^n\to\{-1,1\}$ be a sign-valued Boolean function and let
$\delta\in\mathbb R$. The switching-failure probability at parameter $\delta$ and
threshold $0$ is $1$.
\end{theorem}

\begin{theorem}[Nonnegativity of switching failure]
\lean{FABL.switchingFailureProbability_nonneg}
\label{FABL.switchingFailureProbability_nonneg}
\leanok
\uses{FABL.BooleanFunction, FABL.deltaRandomSubsetWeight_nonneg, FABL.expectRandomRestriction, FABL.switchingFailureIndicator_nonneg, FABL.switchingFailureProbability}
\difficulty{3}
Let $f:\{-1,1\}^n\to\{-1,1\}$ be a sign-valued Boolean function, let $k\in\mathbb
N$, and suppose $0\le\delta\le 1$. The switching-failure probability at parameter
$\delta$ and threshold $k$ is nonnegative.
\end{theorem}

\begin{theorem}[Upper bound for switching failure]
\lean{FABL.switchingFailureProbability_le_one}
\label{FABL.switchingFailureProbability_le_one}
\leanok
\uses{FABL.BooleanFunction, FABL.FixedSignCube, FABL.deltaRandomSubsetWeight, FABL.deltaRandomSubsetWeight_nonneg, FABL.expectRandomRestriction, FABL.sum_deltaRandomSubsetWeight, FABL.switchingFailureIndicator, FABL.switchingFailureIndicator_le_one, FABL.switchingFailureProbability}
\difficulty{4}
Let $f:\{-1,1\}^n\to\{-1,1\}$ be a sign-valued Boolean function, let $k\in\mathbb
N$, and suppose $0\le\delta\le 1$. The switching-failure probability, the weighted
average over random restrictions of the indicator that the restricted function has
decision-tree depth at least $k$, is at most $1$.
\end{theorem}

\begin{definition}[Coordinate restriction weight]
\lean{FABL.coordRestrictionWeightAt}
\label{FABL.coordRestrictionWeightAt}
\leanok
\uses{FABL.CoordRestriction}
For a real number $\delta$, define the weight of a coordinate restriction to be $\delta$
when the coordinate is free, and $(1-\delta)/2$ when it is fixed to either $1$ or $-1$.
\end{definition}

\begin{theorem}[Total coordinate restriction weight]
\lean{FABL.sum_coordRestrictionWeightAt}
\label{FABL.sum_coordRestrictionWeightAt}
\leanok
\uses{FABL.CoordRestriction, FABL.coordRestrictionWeightAt}
\difficulty{4}
For a real number $\delta$, assign weight $\delta$ to the free restriction and weight
$(1-\delta)/2$ to each of the two restrictions fixing the coordinate to $1$ or $-1$. The
sum of these three weights is $1$.
\end{theorem}

\begin{theorem}[Nonnegative coordinate weights]
\lean{FABL.coordRestrictionWeightAt_nonneg}
\label{FABL.coordRestrictionWeightAt_nonneg}
\leanok
\uses{FABL.CoordRestriction, FABL.coordRestrictionWeightAt}
\difficulty{4}
If $0\leq\delta\leq1$, then the coordinate weight of every coordinate restriction is
nonnegative.
\end{theorem}

\begin{definition}[Weight of a restriction assignment]
\lean{FABL.restrictionAssignmentWeightAt}
\label{FABL.restrictionAssignmentWeightAt}
\leanok
\uses{FABL.CoordRestriction, FABL.coordRestrictionWeightAt}
For a real parameter $\delta$ and an assignment of one of the states free, fixed at $1$,
or fixed at $-1$ to each of the $n$ coordinates, define its weight to be the product
over all coordinates of the corresponding coordinate restriction weights at $\delta$.
\end{definition}

\begin{theorem}[Nonnegativity of assignment weight]
\lean{FABL.restrictionAssignmentWeightAt_nonneg}
\label{FABL.restrictionAssignmentWeightAt_nonneg}
\leanok
\uses{FABL.CoordRestriction, FABL.coordRestrictionWeightAt_nonneg, FABL.restrictionAssignmentWeightAt}
\difficulty{1}
Let $n$ be a natural number, let $b4$ be a real number with $0\leq b4\leq 1$, and
assign to each of the $n$ coordinates one of the states “free,” “fixed to $1$,” or
“fixed to $-1$.” The product of the coordinate weights, which are $b4$ for a free state
and $(1-b4)/2$ for either fixed state, is nonnegative.
\end{theorem}

\begin{theorem}[Parameterized total restriction weight]
\lean{FABL.sum_restrictionAssignmentWeightAt}
\label{FABL.sum_restrictionAssignmentWeightAt}
\leanok
\uses{FABL.CoordRestriction, FABL.coordRestrictionWeightAt, FABL.restrictionAssignmentWeightAt, FABL.sum_coordRestrictionWeightAt}
\difficulty{3}
For any nonnegative integer $n$ and any real number $\delta$, give each coordinate
restriction the weight $\delta$ if it is free and $(1-\delta)/2$ if it is fixed to
either $1$ or $-1$. Give an assignment of restrictions to the $n$ coordinates the
product of its coordinate weights. The sum over all such assignments is $1$.
\end{theorem}

\begin{definition}[Free coordinates of a restriction]
\lean{FABL.CoordRestriction.freeCoordinates}
\statementsource{lmn-constant-depth-fourier}{pdf-63759bd909d5-p005-b001}
\statementsource{arora-barak-ch13-circuit-lowerbounds}{pdf-52929dbaa962-p002-b003}
\statementsource{Analysis-of-Boolean-Functions-by-Ryan-ODonnell}{
  pdf-ffbe9c7b3306-p075-b002,
  pdf-ffbe9c7b3306-p098-b003
}
\label{FABL.CoordRestriction.freeCoordinates}
\leanok
\uses{FABL.CoordRestriction}
Let $n$ be a nonnegative integer and assign to each of the $n$ coordinates one of three
statuses: free, fixed to $1$, or fixed to $-1$. The free-coordinate set consists of
precisely those coordinates assigned the free status.
\end{definition}

\begin{theorem}[Membership in the free coordinate set]
\lean{FABL.CoordRestriction.mem_freeCoordinates_iff}
\label{FABL.CoordRestriction.mem_freeCoordinates_iff}
\leanok
\uses{FABL.CoordRestriction, FABL.CoordRestriction.freeCoordinates}
\difficulty{1}
Let $n$ be a natural number, let $\rho$ assign a restriction state to each of the $n$
coordinates, and let $i$ be a coordinate. Then $i$ belongs to the set of free
coordinates of $\rho$ if and only if $\rho(i)$ is free.
\end{theorem}

\begin{definition}[Completion of a coordinate restriction]
\lean{FABL.CoordRestriction.complete}
\statementsource{lmn-constant-depth-fourier}{pdf-63759bd909d5-p005-b001}
\statementsource{Analysis-of-Boolean-Functions-by-Ryan-ODonnell}{pdf-ffbe9c7b3306-p098-b004}
\label{FABL.CoordRestriction.complete}
\leanok
\uses{FABL.CoordRestriction, FABL.signEncode}
Let $n$ be a nonnegative integer, assign each coordinate one of the statuses free, fixed
to $1$, or fixed to $-1$, and let $x\in\mathbb F_2^n$. The completed sign vector has,
at each free coordinate, the sign encoding of the corresponding bit of $x$; at
coordinates fixed to $1$ or $-1$, it has that prescribed sign.
\end{definition}

\begin{definition}[Restriction of a Boolean function]
\lean{FABL.CoordRestriction.restrict}
\statementsource{lmn-constant-depth-fourier}{pdf-63759bd909d5-p005-b001}
\statementsource{arora-barak-ch13-circuit-lowerbounds}{pdf-52929dbaa962-p002-b003}
\statementsource{Analysis-of-Boolean-Functions-by-Ryan-ODonnell}{
  pdf-ffbe9c7b3306-p075-b002,
  pdf-ffbe9c7b3306-p098-b004
}
\label{FABL.CoordRestriction.restrict}
\leanok
\uses{FABL.BooleanFunction, FABL.CoordRestriction, FABL.CoordRestriction.complete, FABL.Sign}
Let $n$ be a nonnegative integer, let $f$ be a function from the sign cube
$\{−1,1\}^n$ to $\{−1,1\}$, and assign each coordinate one of the statuses free,
fixed to $1$, or fixed to $-1$. The restricted function on $\mathbb F_2^n$ evaluates
$f$ at the sign vector obtained by completing that coordinate restriction using the
input bits.
\end{definition}

\begin{definition}[Coordinate switching failure indicator]
\lean{FABL.coordSwitchingFailureIndicator}
\statementsource{lmn-constant-depth-fourier}{
  pdf-63759bd909d5-p005-b001,
  pdf-63759bd909d5-p005-b004
}
\statementsource{arora-barak-ch13-circuit-lowerbounds}{pdf-52929dbaa962-p002-b003}
\statementsource{Analysis-of-Boolean-Functions-by-Ryan-ODonnell}{pdf-ffbe9c7b3306-p100-b004}
\label{FABL.coordSwitchingFailureIndicator}
\leanok
\uses{FABL.BooleanFunction, FABL.CoordRestriction, FABL.CoordRestriction.restrict, FABL.F₂DecisionTree.decisionTreeDepth}
Let $n$ be a nonnegative integer, let $f$ be a Boolean function on $n$ sign coordinates,
let $k$ be a nonnegative integer, and assign each coordinate one of the statuses free,
fixed to $1$, or fixed to $-1$. The coordinate switching failure indicator is $1$ if the
decision-tree depth of the resulting restricted function is at least $k$, and is $0$
otherwise.
\end{definition}

\begin{definition}[Coordinate switching failure probability]
\lean{FABL.coordSwitchingFailureProbability}
\statementsource{lmn-constant-depth-fourier}{
  pdf-63759bd909d5-p005-b001,
  pdf-63759bd909d5-p005-b004
}
\statementsource{arora-barak-ch13-circuit-lowerbounds}{pdf-52929dbaa962-p002-b003}
\statementsource{Analysis-of-Boolean-Functions-by-Ryan-ODonnell}{pdf-ffbe9c7b3306-p100-b004}
\label{FABL.coordSwitchingFailureProbability}
\leanok
\uses{FABL.BooleanFunction, FABL.CoordRestriction, FABL.coordSwitchingFailureIndicator, FABL.restrictionAssignmentWeightAt}
Let $n$ be a nonnegative integer, let $f$ be a Boolean function on $n$ sign coordinates,
let $\delta\in\mathbb R$, and let $k$ be a nonnegative integer. The coordinate
switching failure probability is the sum over all assignments of the three restriction
statuses to the $n$ coordinates, with each assignment weighted by its
restriction-assignment weight at parameter $\delta$, of the corresponding coordinate
switching failure indicator at threshold $k$.
\end{definition}

\begin{theorem}[Nonnegativity of the failure indicator]
\lean{FABL.coordSwitchingFailureIndicator_nonneg}
\label{FABL.coordSwitchingFailureIndicator_nonneg}
\leanok
\uses{FABL.BooleanFunction, FABL.CoordRestriction, FABL.coordSwitchingFailureIndicator}
\difficulty{4}
For a Boolean function $f:\{-1,1\}^n\to\{\pm1\}$, a nonnegative integer $k$, and a
coordinate restriction $\rho:[n]\to\{\text{free},\text{fixed }1,\text{fixed
}-1\}$, the coordinate-switching failure indicator is nonnegative. Here the indicator
equals $1$ when the decision-tree depth of the function obtained by applying $\rho$ is
at least $k$, and equals $0$ otherwise.
\end{theorem}

\begin{theorem}[Nonnegativity of failure probability]
\lean{FABL.coordSwitchingFailureProbability_nonneg}
\label{FABL.coordSwitchingFailureProbability_nonneg}
\leanok
\uses{FABL.BooleanFunction, FABL.coordSwitchingFailureIndicator_nonneg, FABL.coordSwitchingFailureProbability, FABL.restrictionAssignmentWeightAt_nonneg}
\difficulty{3}
Let $f:\{-1,1\}^n\to\{\pm1\}$ be a Boolean function, let $k$ be a natural number,
and let $\delta\in\mathbb R$ satisfy $0\le\delta\le1$. Then the
coordinate-switching failure probability of $f$ at parameter $\delta$ and threshold $k$
is nonnegative.
\end{theorem}

\begin{theorem}[Upper bound on failure probability]
\lean{FABL.coordSwitchingFailureProbability_le_one}
\label{FABL.coordSwitchingFailureProbability_le_one}
\leanok
\uses{FABL.BooleanFunction, FABL.CoordRestriction, FABL.coordSwitchingFailureIndicator, FABL.coordSwitchingFailureProbability, FABL.restrictionAssignmentWeightAt, FABL.restrictionAssignmentWeightAt_nonneg, FABL.sum_restrictionAssignmentWeightAt}
\difficulty{4}
Let $f:\{-1,1\}^n\to\{\pm1\}$ be a Boolean function, let $k$ be a natural number,
and let $\delta\in\mathbb R$ satisfy $0\le\delta\le1$. Then the
coordinate-switching failure probability of $f$ at parameter $\delta$ and threshold $k$
is at most $1$.
\end{theorem}

\begin{definition}[Coordinate restriction induced by fixed signs]
\lean{FABL.coordRestrictionOf}
\statementsource{lmn-constant-depth-fourier}{pdf-63759bd909d5-p005-b001}
\statementsource{arora-barak-ch13-circuit-lowerbounds}{pdf-52929dbaa962-p002-b003}
\statementsource{Analysis-of-Boolean-Functions-by-Ryan-ODonnell}{
  pdf-ffbe9c7b3306-p075-b002,
  pdf-ffbe9c7b3306-p098-b003
}
\label{FABL.coordRestrictionOf}
\leanok
\uses{FABL.CoordRestriction, FABL.FixedSignCube, FABL.Sign}
Let $n$ be a natural number, let $J\subseteq\{0,\ldots,n-1\}$, and assign a sign
$z_i\in\{\pm1\}$ to each index outside $J$. Define a coordinate restriction on
$\{0,\ldots,n-1\}$ by leaving every coordinate in $J$ free, fixing each coordinate
outside $J$ to $+1$ when its assigned sign is $+1$, and fixing it to $-1$ otherwise.
\end{definition}

\begin{theorem}[Free coordinates of a restriction]
\lean{FABL.coordRestrictionOf_apply_of_mem}
\label{FABL.coordRestrictionOf_apply_of_mem}
\leanok
\uses{FABL.FixedSignCube, FABL.coordRestrictionOf}
\difficulty{1}
Let $n$ be a natural number, let $J\subseteq\{0,\ldots,n-1\}$, and assign a sign in
$\{1,-1\}$ to each coordinate outside $J$. The coordinate restriction determined by
these data is free at every coordinate in $J$.
\end{theorem}

\begin{theorem}[Restriction at a fixed coordinate]
\lean{FABL.coordRestrictionOf_apply_of_not_mem}
\label{FABL.coordRestrictionOf_apply_of_not_mem}
\leanok
\uses{FABL.FixedSignCube, FABL.Sign, FABL.coordRestrictionOf}
\difficulty{1}
Let $n$ be a natural number, let $J\subseteq\{0,\ldots,n-1\}$, and assign a sign in
$\{1,-1\}$ to each coordinate outside $J$. At any coordinate $i\notin J$, the
determined coordinate restriction fixes the coordinate to $1$ when its assigned sign is
$1$, and to $-1$ otherwise.
\end{theorem}

\begin{theorem}[Free coordinates of a coordinate restriction]
\lean{FABL.freeCoordinates_coordRestrictionOf}
\label{FABL.freeCoordinates_coordRestrictionOf}
\leanok
\uses{FABL.CoordRestriction.freeCoordinates, FABL.FixedSignCube, FABL.coordRestrictionOf, FABL.coordRestrictionOf_apply_of_not_mem}
\difficulty{4}
Let $n$ be a nonnegative integer, let $J\subseteq\{0,\ldots,n-1\}$, and assign a sign in
$\{-1,1\}$ to each coordinate outside $J$. Form a coordinate restriction by leaving the
coordinates in $J$ free and fixing each other coordinate to its assigned sign. The set
of free coordinates of this restriction is exactly $J$.
\end{theorem}

\begin{theorem}[Completion of a coordinate restriction]
\lean{FABL.complete_coordRestrictionOf}
\label{FABL.complete_coordRestrictionOf}
\leanok
\uses{FABL.CoordRestriction.complete, FABL.FixedSignCube, FABL.combineSignCube, FABL.combineSignCube_apply_fixed, FABL.combineSignCube_apply_free, FABL.coordRestrictionOf, FABL.signEncode}
\difficulty{4}
Let $n$ be a natural number, let $J$ be a subset of $	ext{Fin}(n)$, let $z$ assign a
sign in ${ 1,1}$ to every coordinate outside $J$, and let $x	ext{in }F_2^n$.
Completing the coordinate restriction that leaves $J$ free and fixes the other
coordinates according to $z$ gives the sign vector equal to $x$ encoded as signs on $J$
and to $z$ on the complement of $J$.
\end{theorem}

\begin{theorem}[Binary form of a fixed sign restriction]
\lean{FABL.restrict_coordRestrictionOf}
\label{FABL.restrict_coordRestrictionOf}
\leanok
\uses{FABL.BooleanFunction, FABL.CoordRestriction.restrict, FABL.FixedSignCube, FABL.complete_coordRestrictionOf, FABL.coordRestrictionOf, FABL.extendedSignRestriction, FABL.restrictedBinaryFunction}
\difficulty{3}
Let $f$ be a function from the sign cube on $n$ coordinates to the signs, let $J$ be a
set of coordinates, and fix signs $z$ outside $J$. Applying the coordinate rules that
leave $J$ free and fix every coordinate outside $J$ according to $z$, then converting a
binary input to signs and evaluating $f$, gives the same function as fixing the outside
signs in $f$ and evaluating it on the sign vector obtained from the binary input.
\end{theorem}

\begin{theorem}[Switching failure for fixed coordinates]
\lean{FABL.coordSwitchingFailureIndicator_coordRestrictionOf}
\label{FABL.coordSwitchingFailureIndicator_coordRestrictionOf}
\leanok
\uses{FABL.BooleanFunction, FABL.FixedSignCube, FABL.coordRestrictionOf, FABL.coordSwitchingFailureIndicator, FABL.restrict_coordRestrictionOf, FABL.restrictedDecisionTreeDepth, FABL.switchingFailureIndicator}
\difficulty{1}
Let $f:\{\pm1\}^n\to\{\pm1\}$ be a Boolean function, let $k$ be a natural number, let
$J\subseteq\text{Fin}(n)$, and fix signs $z$ on the coordinates outside $J$. The
switching failure indicator of $f$ at threshold $k$ under the coordinate assignment that
leaves $J$ free and fixes the other coordinates according to $z$ equals the switching
failure indicator at threshold $k$ for the function obtained by fixing those outside
coordinates to $z$.
\end{theorem}

\begin{theorem}[Weight of a coordinate restriction]
\lean{FABL.restrictionAssignmentWeightAt_coordRestrictionOf}
\label{FABL.restrictionAssignmentWeightAt_coordRestrictionOf}
\leanok
\uses{FABL.FixedSignCube, FABL.coordRestrictionOf, FABL.coordRestrictionOf_apply_of_not_mem, FABL.coordRestrictionWeightAt, FABL.restrictionAssignmentWeightAt}
\difficulty{4}
Let $n$ be a natural number, let $\delta\in\mathbb R$, let $J$ be a set of coordinates
in $[n]$, and assign a sign in $\{\pm1\}$ to each coordinate outside $J$. The weight of
the resulting restriction, which leaves coordinates in $J$ free and fixes the others to
their assigned signs, is $\delta^{|J|}\bigl((1-\delta)/2\bigr)^{n-|J|}$.
\end{theorem}

\begin{definition}[Signs fixed by a coordinate restriction]
\lean{FABL.CoordRestriction.fixedSign}
\statementsource{Analysis-of-Boolean-Functions-by-Ryan-ODonnell}{pdf-ffbe9c7b3306-p098-b003}
\label{FABL.CoordRestriction.fixedSign}
\leanok
\uses{FABL.CoordRestriction, FABL.CoordRestriction.freeCoordinates, FABL.FixedSignCube}
Let $n$ be a natural number and let $\rho$ assign to each of the $n$ coordinates one of
three statuses: free, fixed to $+1$, or fixed to $-1$. On the coordinates that are not
free, define the fixed sign assignment to be $+1$ where $\rho$ fixes the coordinate to
$+1$, and $-1$ everywhere else.
\end{definition}

\begin{theorem}[Recovery of a coordinate restriction]
\lean{FABL.coordRestrictionOf_freeCoordinates_fixedSign}
\label{FABL.coordRestrictionOf_freeCoordinates_fixedSign}
\leanok
\uses{FABL.CoordRestriction, FABL.CoordRestriction.fixedSign, FABL.CoordRestriction.freeCoordinates, FABL.CoordRestriction.mem_freeCoordinates_iff, FABL.coordRestrictionOf}
\difficulty{4}
Let $n$ be a natural number and let $\rho$ assign to each of the $n$ coordinates one of
three states: free, fixed to $1$, or fixed to $-1$. If $J$ is the set of free
coordinates and each coordinate outside $J$ is assigned the sign specified by $\rho$,
then the restriction determined by $J$ and those signs is exactly $\rho$.
\end{theorem}

\begin{definition}[Random restriction atom]
\lean{FABL.RandomRestrictionAtom}
\statementsource{lmn-constant-depth-fourier}{pdf-63759bd909d5-p005-b001}
\statementsource{arora-barak-ch13-circuit-lowerbounds}{pdf-52929dbaa962-p002-b003}
\label{FABL.RandomRestrictionAtom}
\leanok
\uses{FABL.FixedSignCube}
For a natural number $n$, a random restriction atom consists of a subset $J$ of the $n$
coordinates together with an assignment of a sign in $\{\pm1\}$ to each coordinate
outside $J$.
\end{definition}

\begin{definition}[Restriction atom correspondence]
\lean{FABL.randomRestrictionAtomEquiv}
\label{FABL.randomRestrictionAtomEquiv}
\leanok
\uses{FABL.CoordRestriction, FABL.CoordRestriction.fixedSign, FABL.CoordRestriction.freeCoordinates, FABL.RandomRestrictionAtom, FABL.coordRestrictionOf, FABL.coordRestrictionOf_freeCoordinates_fixedSign, FABL.freeCoordinates_coordRestrictionOf}
For each $n$, the set of atoms consisting of a subset $J$ of coordinates together with a
sign in $\{\pm1\}$ at every coordinate outside $J$ is in bijection with the assignments
of a restriction state (free, fixed at $1$, or fixed at $-1$) to each coordinate. The
bijection associates to an atom the restriction that is free on $J$ and fixed according
to the atom's sign at every coordinate outside $J$.
\end{definition}

\begin{theorem}[Cardinality of a fixed sign cube]
\lean{FABL.card_fixedSignCube}
\label{FABL.card_fixedSignCube}
\leanok
\uses{FABL.FixedIndex, FABL.FixedSignCube, FABL.Sign}
\difficulty{3}
For a natural number $n$ and a set $J$ of indices among $n$ coordinates, the set of sign
assignments to the coordinates outside $J$ has cardinality $2^{n-|J|}$.
\end{theorem}

\begin{theorem}[Restriction weight and random subset weight]
\lean{FABL.restrictionAssignmentWeightAt_coordRestrictionOf_eq_div}
\label{FABL.restrictionAssignmentWeightAt_coordRestrictionOf_eq_div}
\leanok
\uses{FABL.FixedSignCube, FABL.card_fixedSignCube, FABL.coordRestrictionOf, FABL.deltaRandomSubsetWeight, FABL.restrictionAssignmentWeightAt, FABL.restrictionAssignmentWeightAt_coordRestrictionOf}
\difficulty{3}
Let $n$ be a natural number, let $\delta\in\mathbb R$, let $J$ be a set of coordinates
in $[n]$, and assign a sign in $\{\pm1\}$ to each coordinate outside $J$. The weight of
the restriction that leaves $J$ free and fixes the remaining coordinates to those signs
equals $\delta^{|J|}(1-\delta)^{n-|J|}$ divided by the number of sign assignments on the
coordinates outside $J$.
\end{theorem}

\begin{theorem}[Two forms of failure probability]
\lean{FABL.coordSwitchingFailureProbability_eq}
\label{FABL.coordSwitchingFailureProbability_eq}
\leanok
\uses{FABL.BooleanFunction, FABL.FixedSignCube, FABL.RandomRestrictionAtom, FABL.coordSwitchingFailureIndicator, FABL.coordSwitchingFailureIndicator_coordRestrictionOf, FABL.coordSwitchingFailureProbability, FABL.deltaRandomSubsetWeight, FABL.expectRandomRestriction, FABL.randomRestrictionAtomEquiv, FABL.restrictionAssignmentWeightAt, FABL.restrictionAssignmentWeightAt_coordRestrictionOf_eq_div, FABL.switchingFailureIndicator, FABL.switchingFailureProbability}
\difficulty{4}
Let $f:\{-1,1\}^n\to\{\pm1\}$ be a Boolean function, let $\delta\in\mathbb R$,
and let $k$ be a natural number. The coordinate-switching failure probability, summed
over coordinate restrictions with each free coordinate assigned weight $\delta$ and
each of the two fixed states assigned weight $(1-\delta)/2$, equals the expected
failure indicator obtained by first choosing a free-coordinate set with weight
$\delta^{|J|}(1-\delta)^{n-|J|}$ and then choosing signs on the remaining coordinates
uniformly.
\end{theorem}

\begin{definition}[Compatibility of a term and restriction]
\lean{FABL.DNFTerm.Compatible}
\statementsource{arora-barak-ch13-circuit-lowerbounds}{pdf-52929dbaa962-p003-b004}
\statementsource{Analysis-of-Boolean-Functions-by-Ryan-ODonnell}{pdf-ffbe9c7b3306-p098-b004}
\label{FABL.DNFTerm.Compatible}
\leanok
\uses{FABL.CoordRestriction, FABL.CoordRestriction.complete, FABL.DNFTerm, FABL.DNFTerm.eval}
Let $n$ be a natural number, let $T$ be a conjunction of literals on $n$ coordinates
with distinct coordinates, and let $\rho$ assign each coordinate a status of free, fixed
to $+1$, or fixed to $-1$. The term is compatible with $\rho$ if there exists a binary
assignment to the free coordinates such that, after completing it using the fixed values
specified by $\rho$, the term evaluates to $-1$.
\end{definition}

\begin{definition}[Bad restriction for a DNF formula]
\lean{FABL.DNFFormula.IsBadRestriction}
\statementsource{lmn-constant-depth-fourier}{pdf-63759bd909d5-p005-b004}
\statementsource{arora-barak-ch13-circuit-lowerbounds}{pdf-52929dbaa962-p003-b004}
\statementsource{Analysis-of-Boolean-Functions-by-Ryan-ODonnell}{pdf-ffbe9c7b3306-p100-b004}
\label{FABL.DNFFormula.IsBadRestriction}
\leanok
\uses{FABL.CoordRestriction, FABL.CoordRestriction.restrict, FABL.DNFFormula, FABL.DNFFormula.toBooleanFunction}
Let $n$ be a natural number, let $\varphi$ be a disjunction of conjunctions of literals
on $n$ coordinates, and let $\rho$ assign each coordinate a status of free, fixed to
$+1$, or fixed to $-1$. The restriction is bad for $\varphi$ if there are two binary
assignments to the free coordinates for which the restricted formula has different
values.
\end{definition}

\begin{theorem}[A bad restriction admits a compatible term]
\lean{FABL.DNFFormula.exists_compatible_of_badRestriction}
\label{FABL.DNFFormula.exists_compatible_of_badRestriction}
\leanok
\uses{FABL.CoordRestriction, FABL.CoordRestriction.complete, FABL.CoordRestriction.restrict, FABL.DNFFormula, FABL.DNFFormula.IsBadRestriction, FABL.DNFFormula.eval, FABL.DNFFormula.eval_eq_neg_one_iff, FABL.DNFFormula.toBooleanFunction, FABL.DNFTerm.Compatible}
\difficulty{5}
Let $n$ be a nonnegative integer, let $\varphi$ be a DNF formula on $n$ sign
coordinates, and let $\rho$ be a coordinate restriction. If the restricted Boolean
function is nonconstant, then some term of $\varphi$ is compatible with $\rho$, meaning
it evaluates to $-1$ on at least one completion of that restriction.
\end{theorem}

\begin{theorem}[A free literal under a bad restriction]
\lean{FABL.DNFTerm.exists_free_literal_of_compatible_of_badRestriction}
\proofsource{arora-barak-ch13-circuit-lowerbounds}{
  pdf-52929dbaa962-p003-b005,
  pdf-52929dbaa962-p004-b001
}
\proofsource{Analysis-of-Boolean-Functions-by-Ryan-ODonnell}{pdf-ffbe9c7b3306-p110-b003}
\label{FABL.DNFTerm.exists_free_literal_of_compatible_of_badRestriction}
\leanok
\uses{FABL.CoordRestriction, FABL.CoordRestriction.complete, FABL.CoordRestriction.restrict, FABL.DNFFormula, FABL.DNFFormula.IsBadRestriction, FABL.DNFFormula.eval, FABL.DNFFormula.eval_eq_neg_one_iff, FABL.DNFFormula.toBooleanFunction, FABL.DNFTerm, FABL.DNFTerm.Compatible, FABL.DNFTerm.eval, FABL.DNFTerm.eval_eq_neg_one_iff}
\difficulty{5}
Let $\varphi$ be a disjunction of conjunctions of literals on $n$ coordinates, let $T$
be one of its terms, and let $\rho$ assign each coordinate either free, fixed to $1$, or
fixed to $-1$. Suppose $T$ is compatible with $\rho$, meaning some assignment to the
free coordinates makes every literal of $T$ hold after applying the restriction. Suppose
also that the restricted function defined by $\varphi$ takes different values on two
assignments to the free coordinates. Then $T$ contains a literal whose coordinate is
free under $\rho$.
\end{theorem}

\begin{definition}[Indices of compatible terms]
\lean{FABL.DNFFormula.compatibleTermIndices}
\statementsource{arora-barak-ch13-circuit-lowerbounds}{pdf-52929dbaa962-p003-b005}
\label{FABL.DNFFormula.compatibleTermIndices}
\leanok
\uses{FABL.CoordRestriction, FABL.DNFFormula, FABL.DNFTerm.Compatible}
Let $n$ be a natural number, let $\varphi$ be a finite list of conjunctions of literals
on $n$ coordinates, and let $\rho$ assign each coordinate a status of free, fixed to
$+1$, or fixed to $-1$. The compatible term indices are the indices in the formula's
list whose corresponding term is compatible with $\rho$.
\end{definition}

\begin{theorem}[Compatible term under a bad restriction]
\lean{FABL.DNFFormula.compatibleTermIndices_nonempty_of_badRestriction}
\proofsource{arora-barak-ch13-circuit-lowerbounds}{pdf-52929dbaa962-p003-b005}
\label{FABL.DNFFormula.compatibleTermIndices_nonempty_of_badRestriction}
\leanok
\uses{FABL.CoordRestriction, FABL.DNFFormula, FABL.DNFFormula.IsBadRestriction, FABL.DNFFormula.compatibleTermIndices, FABL.DNFFormula.exists_compatible_of_badRestriction}
\difficulty{4}
Let $\phi$ be a DNF formula on $n$ coordinates and let $\rho$ be a coordinate
restriction for which the restricted Boolean function is nonconstant. Then at least one
term of $\phi$ is compatible with $\rho$.
\end{theorem}

\begin{definition}[First compatible term index]
\lean{FABL.DNFFormula.firstCompatibleTermIndex}
\statementsource{arora-barak-ch13-circuit-lowerbounds}{
  pdf-52929dbaa962-p004-b001,
  pdf-52929dbaa962-p004-b002
}
\label{FABL.DNFFormula.firstCompatibleTermIndex}
\leanok
\uses{FABL.CoordRestriction, FABL.DNFFormula, FABL.DNFFormula.IsBadRestriction, FABL.DNFFormula.compatibleTermIndices, FABL.DNFFormula.compatibleTermIndices_nonempty_of_badRestriction}
Let $n$ be a natural number, let $\varphi$ be a finite DNF formula on $n$ coordinates,
and let $\rho$ be a coordinate restriction for which the restricted formula takes
different values on some two binary assignments to the free coordinates. The first
compatible term index is the least index in the formula's list whose term is compatible
with $\rho$.
\end{definition}

\begin{theorem}[First compatible index is compatible]
\lean{FABL.DNFFormula.firstCompatibleTermIndex_mem}
\label{FABL.DNFFormula.firstCompatibleTermIndex_mem}
\leanok
\uses{FABL.CoordRestriction, FABL.DNFFormula, FABL.DNFFormula.IsBadRestriction, FABL.DNFFormula.compatibleTermIndices, FABL.DNFFormula.firstCompatibleTermIndex}
\difficulty{1}
Let $\varphi$ be a DNF formula on $n$ sign-valued Boolean variables, and let $\rho$ be a
coordinate restriction for which the restricted function is nonconstant. The first
compatible term index belongs to the set of indices of terms compatible with $\rho$.
\end{theorem}

\begin{theorem}[The first compatible term is compatible]
\lean{FABL.DNFFormula.firstCompatibleTerm_compatible}
\label{FABL.DNFFormula.firstCompatibleTerm_compatible}
\leanok
\uses{FABL.CoordRestriction, FABL.DNFFormula, FABL.DNFFormula.IsBadRestriction, FABL.DNFFormula.firstCompatibleTermIndex, FABL.DNFFormula.firstCompatibleTermIndex_mem, FABL.DNFTerm.Compatible}
\difficulty{1}
Let $\varphi$ be a DNF formula on $n$ sign-valued Boolean variables, and let $\rho$ be a
coordinate restriction for which the restricted function is nonconstant. The term at the
first compatible term index is compatible with $\rho$.
\end{theorem}

\begin{definition}[Free support of the first compatible term]
\lean{FABL.DNFFormula.firstCompatibleFreeSupport}
\statementsource{arora-barak-ch13-circuit-lowerbounds}{
  pdf-52929dbaa962-p004-b001,
  pdf-52929dbaa962-p004-b002
}
\label{FABL.DNFFormula.firstCompatibleFreeSupport}
\leanok
\uses{FABL.CoordRestriction, FABL.CoordRestriction.freeCoordinates, FABL.DNFFormula, FABL.DNFFormula.IsBadRestriction, FABL.DNFFormula.firstCompatibleTermIndex, FABL.DNFTerm.support}
Let $n$ be a natural number, let $\varphi$ be a finite DNF formula on $n$ coordinates,
and let $\rho$ be a coordinate restriction for which the restricted formula takes
different values on some two binary assignments to the free coordinates. The free
support of the first compatible term is the set of coordinates that are both free under
$\rho$ and occur in the term at the least index whose term is compatible with $\rho$.
\end{definition}

\begin{theorem}[Nonempty free support of a compatible term]
\lean{FABL.DNFFormula.firstCompatibleFreeSupport_nonempty}
\proofsource{arora-barak-ch13-circuit-lowerbounds}{pdf-52929dbaa962-p004-b002}
\label{FABL.DNFFormula.firstCompatibleFreeSupport_nonempty}
\leanok
\uses{FABL.CoordRestriction, FABL.CoordRestriction.freeCoordinates, FABL.CoordRestriction.mem_freeCoordinates_iff, FABL.DNFFormula, FABL.DNFFormula.IsBadRestriction, FABL.DNFFormula.firstCompatibleFreeSupport, FABL.DNFFormula.firstCompatibleTermIndex, FABL.DNFFormula.firstCompatibleTerm_compatible, FABL.DNFTerm.Compatible, FABL.DNFTerm.exists_free_literal_of_compatible_of_badRestriction, FABL.DNFTerm.mem_support_of_mem_literals, FABL.DNFTerm.support}
\difficulty{4}
Let $\varphi$ be a DNF formula on $n$ sign-valued Boolean variables, and let $\rho$ be a
coordinate restriction for which the restricted function is nonconstant. The support of
the first term compatible with $\rho$ has a nonempty intersection with the coordinates
left free by $\rho$.
\end{theorem}

\begin{definition}[First bad coordinate]
\lean{FABL.DNFFormula.firstBadCoordinate}
\statementsource{arora-barak-ch13-circuit-lowerbounds}{pdf-52929dbaa962-p004-b002}
\label{FABL.DNFFormula.firstBadCoordinate}
\leanok
\uses{FABL.CoordRestriction, FABL.DNFFormula, FABL.DNFFormula.IsBadRestriction, FABL.DNFFormula.firstCompatibleFreeSupport, FABL.DNFFormula.firstCompatibleFreeSupport_nonempty}
Let $n$ be a nonnegative integer, let $\varphi$ be a DNF formula on $\{\pm1\}^n$, and
let $\rho$ be a coordinate restriction for which the restricted formula is nonconstant.
The first bad coordinate is the least coordinate in the nonempty set of free coordinates
compatible with a term witnessing this nonconstancy.
\end{definition}

\begin{theorem}[Selected coordinate belongs to the free support]
\lean{FABL.DNFFormula.firstBadCoordinate_mem}
\label{FABL.DNFFormula.firstBadCoordinate_mem}
\leanok
\uses{FABL.CoordRestriction, FABL.DNFFormula, FABL.DNFFormula.IsBadRestriction, FABL.DNFFormula.firstBadCoordinate, FABL.DNFFormula.firstCompatibleFreeSupport}
\difficulty{1}
Let $\varphi$ be a DNF formula on $n$ sign-valued Boolean variables, and let $\rho$ be a
coordinate restriction for which the restricted function is nonconstant. The first
coordinate in the support of the first term compatible with $\rho$, among coordinates
left free by $\rho$, belongs to that term's support intersected with the set of free
coordinates.
\end{theorem}

\begin{theorem}[The selected coordinate is free]
\lean{FABL.DNFFormula.firstBadCoordinate_free}
\proofsource{Analysis-of-Boolean-Functions-by-Ryan-ODonnell}{pdf-ffbe9c7b3306-p110-b003}
\label{FABL.DNFFormula.firstBadCoordinate_free}
\leanok
\uses{FABL.CoordRestriction, FABL.CoordRestriction.mem_freeCoordinates_iff, FABL.DNFFormula, FABL.DNFFormula.IsBadRestriction, FABL.DNFFormula.firstBadCoordinate, FABL.DNFFormula.firstBadCoordinate_mem}
\difficulty{1}
Let $\varphi$ be a DNF formula on $n$ sign-valued Boolean variables, and let $\rho$ be a
coordinate restriction for which the restricted function is nonconstant. The first
coordinate in the support of the first term compatible with $\rho$, among coordinates
left free by $\rho$, is itself free under $\rho$.
\end{theorem}

\begin{definition}[Fixed state]
\lean{FABL.CoordRestriction.fixedState}
\statementsource{Analysis-of-Boolean-Functions-by-Ryan-ODonnell}{pdf-ffbe9c7b3306-p098-b003}
\label{FABL.CoordRestriction.fixedState}
\leanok
\uses{FABL.CoordRestriction, FABL.Sign}
For a sign $s\in\{\pm1\}$, define its fixed state to be the restriction fixing the
coordinate to $+1$ when $s=+1$, and to $-1$ otherwise.
\end{definition}

\begin{definition}[Fixing a coordinate]
\lean{FABL.CoordRestriction.fixCoordinate}
\statementsource{Analysis-of-Boolean-Functions-by-Ryan-ODonnell}{pdf-ffbe9c7b3306-p110-b003}
\label{FABL.CoordRestriction.fixCoordinate}
\leanok
\uses{FABL.CoordRestriction, FABL.CoordRestriction.fixedState, FABL.Sign}
Let $n$ be a nonnegative integer, let $\rho$ be a coordinate restriction on $n$
coordinates, let $i$ be one of those coordinates, and let $s\in\{\pm1\}$. The
restriction obtained by fixing coordinate $i$ according to $s$ agrees with $\rho$ at
every other coordinate, and fixes coordinate $i$ to $+1$ if $s=+1$ and to $-1$
otherwise.
\end{definition}

\begin{theorem}[State at the fixed coordinate]
\lean{FABL.CoordRestriction.fixCoordinate_apply_self}
\label{FABL.CoordRestriction.fixCoordinate_apply_self}
\leanok
\uses{FABL.CoordRestriction, FABL.CoordRestriction.fixCoordinate, FABL.CoordRestriction.fixedState, FABL.Sign}
\difficulty{1}
Let $n$ be a natural number, let $\rho$ assign one of the three restriction states to
each of the $n$ coordinates, let $i$ be an $n$-coordinate index, and let $s\in\{\pm1\}$.
After fixing coordinate $i$ according to $s$, its restriction state is the state
determined by $s$.
\end{theorem}

\begin{theorem}[Other coordinates remain unchanged]
\lean{FABL.CoordRestriction.fixCoordinate_apply_of_ne}
\label{FABL.CoordRestriction.fixCoordinate_apply_of_ne}
\leanok
\uses{FABL.CoordRestriction, FABL.CoordRestriction.fixCoordinate, FABL.Sign}
\difficulty{1}
Let $n$ be a natural number, let $\rho$ assign one of the three restriction states to
each of the $n$ coordinates, and let $i,j$ be distinct coordinates among the $n$
coordinates. Fixing coordinate $i$ according to a sign $s\in\{\pm1\}$ leaves the
restriction state at coordinate $j$ equal to its original state $\rho(j)$.
\end{theorem}

\begin{theorem}[Completion after fixing a free coordinate]
\lean{FABL.CoordRestriction.complete_fixCoordinate_of_free}
\label{FABL.CoordRestriction.complete_fixCoordinate_of_free}
\leanok
\uses{FABL.CoordRestriction, FABL.CoordRestriction.complete, FABL.CoordRestriction.fixCoordinate, FABL.CoordRestriction.fixedState, FABL.Sign, FABL.binarySignEquiv, FABL.signEncode}
\difficulty{4}
Let $n$ be a natural number, let $\rho$ be a coordinate restriction on $n$ coordinates,
let $i$ be a coordinate that is free in $\rho$, let $s$ be a sign, and let $x$ be a
binary vector of length $n$. Completing the restriction obtained by fixing $i$ to $s$,
using $x$, gives the same sign vector as completing $\rho$ after replacing the $i$th
entry of $x$ by the binary value corresponding to $s$.
\end{theorem}

\begin{theorem}[Completion after fixing a coordinate to its value]
\lean{FABL.CoordRestriction.complete_fixCoordinate_eq_of_value}
\label{FABL.CoordRestriction.complete_fixCoordinate_eq_of_value}
\leanok
\uses{FABL.CoordRestriction, FABL.CoordRestriction.complete, FABL.CoordRestriction.fixCoordinate, FABL.CoordRestriction.fixedState, FABL.Sign}
\difficulty{4}
Let $n$ be a natural number, let $\rho$ be a coordinate restriction on $n$ coordinates,
let $i$ be a coordinate, let $s$ be a sign, and let $x$ be a binary vector of length
$n$. If the completion of $\rho$ using $x$ has value $s$ at coordinate $i$, then
completing the restriction obtained by fixing coordinate $i$ to $s$, still using $x$,
leaves the completed sign vector unchanged.
\end{theorem}

\begin{theorem}[Compatibility after freeing a coordinate]
\lean{FABL.DNFTerm.compatible_of_fixCoordinate}
\label{FABL.DNFTerm.compatible_of_fixCoordinate}
\leanok
\uses{FABL.CoordRestriction, FABL.CoordRestriction.complete_fixCoordinate_of_free, FABL.CoordRestriction.fixCoordinate, FABL.DNFTerm, FABL.DNFTerm.Compatible, FABL.Sign, FABL.binarySignEquiv}
\difficulty{4}
Let $T$ be a conjunction of literals on $n$ coordinates, let $\rho$ assign each
coordinate either no restriction or a fixed sign, and let $i$ be a coordinate and
$s\in\{\pm1\}$. If $\rho$ leaves $i$ unrestricted and there is a binary assignment
making $T$ true after $i$ is fixed to the state determined by $s$, then there is a
binary assignment making $T$ true under $\rho$.
\end{theorem}

\begin{definition}[First compatible term index]
\lean{FABL.DNFFormula.firstCompatibleTermIndexOf}
\label{FABL.DNFFormula.firstCompatibleTermIndexOf}
\leanok
\uses{FABL.CoordRestriction, FABL.DNFFormula, FABL.DNFFormula.compatibleTermIndices}
Let $\varphi$ be a disjunctive normal form on $n$ coordinates, and let $\rho$ assign
each coordinate one of the states free, fixed at $1$, or fixed at $-1$. If at least one
term of $\varphi$ is compatible with $\rho$, the selected index is the least index in
the ordered list of terms whose term is compatible with $\rho$.
\end{definition}

\begin{theorem}[Agreement of two first-term selections]
\lean{FABL.DNFFormula.firstCompatibleTermIndex_eq_firstCompatibleTermIndexOf}
\label{FABL.DNFFormula.firstCompatibleTermIndex_eq_firstCompatibleTermIndexOf}
\leanok
\uses{FABL.CoordRestriction, FABL.DNFFormula, FABL.DNFFormula.IsBadRestriction, FABL.DNFFormula.compatibleTermIndices_nonempty_of_badRestriction, FABL.DNFFormula.firstCompatibleTermIndex, FABL.DNFFormula.firstCompatibleTermIndexOf}
\difficulty{1}
Let $\varphi$ be a DNF formula on $n$ sign-valued Boolean variables, and let $\rho$ be a
coordinate restriction for which the restricted function is nonconstant. The first
compatible term index defined using that nonconstancy agrees with the first compatible
term index defined using the resulting nonempty set of compatible indices.
\end{theorem}

\begin{definition}[First bad required sign]
\lean{FABL.DNFFormula.firstBadRequiredSign}
\statementsource{arora-barak-ch13-circuit-lowerbounds}{pdf-52929dbaa962-p004-b002}
\label{FABL.DNFFormula.firstBadRequiredSign}
\leanok
\uses{FABL.CoordRestriction, FABL.DNFFormula, FABL.DNFFormula.IsBadRestriction, FABL.DNFFormula.firstBadCoordinate, FABL.DNFFormula.firstBadCoordinate_mem, FABL.DNFFormula.firstCompatibleTermIndex, FABL.DNFTerm.requiredAt, FABL.Sign}
Let $n$ be a nonnegative integer, let $\varphi$ be a DNF formula on $\{\pm1\}^n$, and
let $\rho$ be a coordinate restriction for which the restricted formula is nonconstant.
The first bad required sign is the sign required at the first bad coordinate by the
first compatible term.
\end{definition}

\begin{definition}[Bad extension]
\lean{FABL.DNFFormula.badExtension}
\statementsource{arora-barak-ch13-circuit-lowerbounds}{
  pdf-52929dbaa962-p003-b004,
  pdf-52929dbaa962-p004-b002
}
\statementsource{Analysis-of-Boolean-Functions-by-Ryan-ODonnell}{pdf-ffbe9c7b3306-p110-b003}
\label{FABL.DNFFormula.badExtension}
\leanok
\uses{FABL.CoordRestriction, FABL.CoordRestriction.fixCoordinate, FABL.DNFFormula, FABL.DNFFormula.IsBadRestriction, FABL.DNFFormula.firstBadCoordinate, FABL.DNFFormula.firstBadRequiredSign}
Let $n$ be a nonnegative integer, let $\varphi$ be a DNF formula on $\{\pm1\}^n$, and
let $\rho$ be a coordinate restriction for which the restricted formula is nonconstant.
Its bad extension is obtained from $\rho$ by fixing the first bad coordinate to the
first bad required sign.
\end{definition}

\begin{theorem}[First compatible term compatible bad extension]
\lean{FABL.DNFFormula.firstCompatibleTerm_compatible_badExtension}
\statementsource{arora-barak-ch13-circuit-lowerbounds}{
  pdf-52929dbaa962-p004-b002,
  pdf-52929dbaa962-p005-b001
}
\proofsource{Analysis-of-Boolean-Functions-by-Ryan-ODonnell}{pdf-ffbe9c7b3306-p110-b003}
\label{FABL.DNFFormula.firstCompatibleTerm_compatible_badExtension}
\leanok
\uses{FABL.CoordRestriction, FABL.CoordRestriction.complete, FABL.CoordRestriction.complete_fixCoordinate_eq_of_value, FABL.DNFFormula, FABL.DNFFormula.IsBadRestriction, FABL.DNFFormula.badExtension, FABL.DNFFormula.firstBadCoordinate, FABL.DNFFormula.firstBadCoordinate_mem, FABL.DNFFormula.firstBadRequiredSign, FABL.DNFFormula.firstCompatibleTermIndex, FABL.DNFFormula.firstCompatibleTerm_compatible, FABL.DNFTerm.Compatible, FABL.DNFTerm.eval_eq_neg_one_iff_requiredAt, FABL.DNFTerm.requiredAt, FABL.DNFTerm.support}
\difficulty{4}
% dataset-pipeline: stub -- prose pending build_informal_statements.py
Statement pending regeneration.
\end{theorem}

\begin{theorem}[Compatible of bad extension]
\lean{FABL.DNFFormula.compatible_of_badExtension}
\statementsource{arora-barak-ch13-circuit-lowerbounds}{pdf-52929dbaa962-p004-b002}
\proofsource{Analysis-of-Boolean-Functions-by-Ryan-ODonnell}{pdf-ffbe9c7b3306-p110-b003}
\label{FABL.DNFFormula.compatible_of_badExtension}
\leanok
\uses{FABL.CoordRestriction, FABL.DNFFormula, FABL.DNFFormula.IsBadRestriction, FABL.DNFFormula.badExtension, FABL.DNFFormula.firstBadCoordinate, FABL.DNFFormula.firstBadCoordinate_free, FABL.DNFFormula.firstBadRequiredSign, FABL.DNFTerm, FABL.DNFTerm.Compatible, FABL.DNFTerm.compatible_of_fixCoordinate}
\difficulty{1}
% dataset-pipeline: stub -- prose pending build_informal_statements.py
Statement pending regeneration.
\end{theorem}

\begin{theorem}[Compatible term after a bad extension]
\lean{FABL.DNFFormula.compatibleTermIndices_badExtension_nonempty}
\proofsource{arora-barak-ch13-circuit-lowerbounds}{pdf-52929dbaa962-p004-b002}
\proofsource{Analysis-of-Boolean-Functions-by-Ryan-ODonnell}{pdf-ffbe9c7b3306-p110-b003}
\label{FABL.DNFFormula.compatibleTermIndices_badExtension_nonempty}
\leanok
\uses{FABL.CoordRestriction, FABL.DNFFormula, FABL.DNFFormula.IsBadRestriction, FABL.DNFFormula.badExtension, FABL.DNFFormula.compatibleTermIndices, FABL.DNFFormula.firstCompatibleTermIndex, FABL.DNFFormula.firstCompatibleTerm_compatible_badExtension}
\difficulty{3}
Let $\phi$ be a DNF formula on $n$ coordinates and let $\rho$ be a coordinate
restriction whose restricted Boolean function is nonconstant. Fix a free coordinate
occurring in the first compatible term and set it to that term's required sign,
obtaining the bad extension of $\rho$. Then at least one term of $\phi$ is compatible
with this extended restriction.
\end{theorem}

\begin{theorem}[First compatible term under a bad extension]
\lean{FABL.DNFFormula.firstCompatibleTermIndexOf_badExtension}
\proofsource{Analysis-of-Boolean-Functions-by-Ryan-ODonnell}{pdf-ffbe9c7b3306-p110-b003}
\label{FABL.DNFFormula.firstCompatibleTermIndexOf_badExtension}
\leanok
\uses{FABL.CoordRestriction, FABL.DNFFormula, FABL.DNFFormula.IsBadRestriction, FABL.DNFFormula.badExtension, FABL.DNFFormula.compatibleTermIndices, FABL.DNFFormula.compatibleTermIndices_badExtension_nonempty, FABL.DNFFormula.compatible_of_badExtension, FABL.DNFFormula.firstCompatibleTermIndex, FABL.DNFFormula.firstCompatibleTermIndexOf, FABL.DNFFormula.firstCompatibleTerm_compatible_badExtension}
\difficulty{4}
Let $\varphi$ be a DNF formula on $n$ sign-valued Boolean variables, and let $\rho$ be a
coordinate restriction for which the restricted function is nonconstant. The first term
compatible with the extension that fixes the selected coordinate to the sign required by
the first compatible term has the same index as the first term compatible with $\rho$.
\end{theorem}

\begin{definition}[Bad restrictions]
\lean{FABL.DNFFormula.badRestrictions}
\statementsource{arora-barak-ch13-circuit-lowerbounds}{pdf-52929dbaa962-p003-b004}
\statementsource{Analysis-of-Boolean-Functions-by-Ryan-ODonnell}{pdf-ffbe9c7b3306-p100-b004}
\label{FABL.DNFFormula.badRestrictions}
\leanok
\uses{FABL.CoordRestriction, FABL.DNFFormula, FABL.DNFFormula.IsBadRestriction}
Let $n$ be a nonnegative integer and let $\varphi$ be a DNF formula on $\{\pm1\}^n$. Its
bad restrictions are all coordinate restrictions for which the restricted formula is
nonconstant.
\end{definition}

\begin{definition}[Bad extension map]
\lean{FABL.DNFFormula.badExtensionMap}
\statementsource{arora-barak-ch13-circuit-lowerbounds}{
  pdf-52929dbaa962-p003-b004,
  pdf-52929dbaa962-p003-b005,
  pdf-52929dbaa962-p004-b001,
  pdf-52929dbaa962-p004-b002,
  pdf-52929dbaa962-p005-b001
}
\statementsource{Analysis-of-Boolean-Functions-by-Ryan-ODonnell}{pdf-ffbe9c7b3306-p110-b003}
\label{FABL.DNFFormula.badExtensionMap}
\leanok
\uses{FABL.CoordRestriction, FABL.DNFFormula, FABL.DNFFormula.IsBadRestriction, FABL.DNFFormula.badExtension}
Let $n$ be a nonnegative integer and let $\varphi$ be a DNF formula on $\{\pm1\}^n$. The
bad extension map sends each coordinate restriction for which the restricted formula is
nonconstant to the restriction obtained by fixing its first bad coordinate to its first
bad required sign.
\end{definition}

\begin{definition}[Bad coordinate map]
\lean{FABL.DNFFormula.badCoordinateMap}
\statementsource{arora-barak-ch13-circuit-lowerbounds}{
  pdf-52929dbaa962-p004-b001,
  pdf-52929dbaa962-p004-b002,
  pdf-52929dbaa962-p005-b001
}
\statementsource{Analysis-of-Boolean-Functions-by-Ryan-ODonnell}{pdf-ffbe9c7b3306-p110-b003}
\label{FABL.DNFFormula.badCoordinateMap}
\leanok
\uses{FABL.CoordRestriction, FABL.DNFFormula, FABL.DNFFormula.IsBadRestriction, FABL.DNFFormula.firstBadCoordinate}
Let $n$ be a nonnegative integer and let $\varphi$ be a DNF formula on $\{\pm1\}^n$. The
bad coordinate map sends each coordinate restriction for which the restricted formula is
nonconstant to its first bad coordinate.
\end{definition}

\begin{theorem}[Recovering a bad coordinate restriction]
\lean{FABL.DNFFormula.recover_badRestriction}
\proofsource{Analysis-of-Boolean-Functions-by-Ryan-ODonnell}{pdf-ffbe9c7b3306-p110-b003}
\label{FABL.DNFFormula.recover_badRestriction}
\leanok
\uses{FABL.CoordRestriction, FABL.CoordRestriction.fixCoordinate, FABL.DNFFormula, FABL.DNFFormula.IsBadRestriction, FABL.DNFFormula.badCoordinateMap, FABL.DNFFormula.badExtension, FABL.DNFFormula.badExtensionMap, FABL.DNFFormula.firstBadCoordinate, FABL.DNFFormula.firstBadCoordinate_free}
\difficulty{4}
Let $n$ be a natural number and let $\varphi$ be a DNF formula on $n$ coordinates. A
coordinate restriction $\rho$ is bad when the restricted Boolean function is
nonconstant. For every such bad restriction, extend it by fixing its selected bad
coordinate to the sign chosen by the formula's construction; changing that coordinate
back to free recovers $\rho$ exactly.
\end{theorem}

\begin{theorem}[Bad coordinate map injective on fiber]
\lean{FABL.DNFFormula.badCoordinateMap_injective_on_fiber}
\statementsource{arora-barak-ch13-circuit-lowerbounds}{
  pdf-52929dbaa962-p004-b002,
  pdf-52929dbaa962-p005-b001
}
\proofsource{Analysis-of-Boolean-Functions-by-Ryan-ODonnell}{pdf-ffbe9c7b3306-p110-b003}
\label{FABL.DNFFormula.badCoordinateMap_injective_on_fiber}
\leanok
\uses{FABL.CoordRestriction, FABL.DNFFormula, FABL.DNFFormula.badCoordinateMap, FABL.DNFFormula.badExtensionMap, FABL.DNFFormula.recover_badRestriction}
\difficulty{3}
% dataset-pipeline: stub -- prose pending build_informal_statements.py
Statement pending regeneration.
\end{theorem}

\begin{definition}[Bad extension fiber]
\lean{FABL.DNFFormula.badExtensionFiber}
\statementsource{arora-barak-ch13-circuit-lowerbounds}{pdf-52929dbaa962-p003-b004}
\statementsource{Analysis-of-Boolean-Functions-by-Ryan-ODonnell}{pdf-ffbe9c7b3306-p110-b003}
\label{FABL.DNFFormula.badExtensionFiber}
\leanok
\uses{FABL.CoordRestriction, FABL.DNFFormula, FABL.DNFFormula.IsBadRestriction, FABL.DNFFormula.badExtensionMap}
Let $\varphi$ be a DNF formula on $n$ coordinates and let $\eta$ be a coordinate
restriction. The bad extension fiber consists of all bad restrictions $\rho$ whose image
under the bad-extension map of $\varphi$ is $\eta$.
\end{definition}

\begin{theorem}[Bad extension fiber bound]
\lean{FABL.DNFFormula.badExtensionFiber_card_le_width}
\proofsource{Analysis-of-Boolean-Functions-by-Ryan-ODonnell}{pdf-ffbe9c7b3306-p110-b003}
\label{FABL.DNFFormula.badExtensionFiber_card_le_width}
\leanok
\uses{FABL.CoordRestriction, FABL.DNFFormula, FABL.DNFFormula.badCoordinateMap, FABL.DNFFormula.badCoordinateMap_injective_on_fiber, FABL.DNFFormula.badExtension, FABL.DNFFormula.badExtensionFiber, FABL.DNFFormula.badExtensionMap, FABL.DNFFormula.compatibleTermIndices, FABL.DNFFormula.compatible_of_badExtension, FABL.DNFFormula.firstBadCoordinate_mem, FABL.DNFFormula.firstCompatibleTermIndex, FABL.DNFFormula.firstCompatibleTerm_compatible_badExtension, FABL.DNFFormula.width, FABL.DNFFormula.width_le_of_mem, FABL.DNFTerm.card_support, FABL.DNFTerm.support, FABL.DNFTerm.width}
\difficulty{5}
Let $n$ be a natural number, let $\varphi$ be a disjunctive normal form on $n$ sign
variables, and let $w$ be a natural number such that the width of $\varphi$ is at most
$w$. For every coordinate restriction $\eta$, the number of bad restrictions whose
designated bad extension is $\eta$ is at most $w$.
\end{theorem}

\begin{theorem}[Weight of a fixed sign state]
\lean{FABL.coordRestrictionWeightAt_fixedState}
\label{FABL.coordRestrictionWeightAt_fixedState}
\leanok
\uses{FABL.CoordRestriction.fixedState, FABL.Sign, FABL.coordRestrictionWeightAt}
\difficulty{4}
For every real number $\delta$ and every sign $s\in\{\pm1\}$, the coordinate weight of
the restriction that fixes the coordinate to $s$ is $(1-\delta)/2$.
\end{theorem}

\begin{theorem}[Weight change under a bad extension]
\lean{FABL.DNFFormula.restrictionAssignmentWeightAt_badExtension_eq}
\proofsource{Analysis-of-Boolean-Functions-by-Ryan-ODonnell}{pdf-ffbe9c7b3306-p110-b003}
\label{FABL.DNFFormula.restrictionAssignmentWeightAt_badExtension_eq}
\leanok
\uses{FABL.CoordRestriction, FABL.CoordRestriction.fixCoordinate_apply_of_ne, FABL.CoordRestriction.fixedState, FABL.DNFFormula, FABL.DNFFormula.IsBadRestriction, FABL.DNFFormula.badExtension, FABL.DNFFormula.firstBadCoordinate, FABL.DNFFormula.firstBadCoordinate_free, FABL.DNFFormula.firstBadRequiredSign, FABL.coordRestrictionWeightAt, FABL.coordRestrictionWeightAt_fixedState, FABL.restrictionAssignmentWeightAt}
\difficulty{5}
Let $n$ be a nonnegative integer, let $\varphi$ be a disjunctive normal form on $n$
coordinates, let $\delta\in\mathbb R$ with $\delta\ne1$, and let $\rho$ be a coordinate
restriction for which the restricted formula is nonconstant. If $\rho'$ is obtained by
fixing the first free coordinate in the support of the first term compatible with $\rho$
to that term's required sign, then the product weight of $\rho$ is
$\frac{2\delta}{1-\delta}$ times the product weight of $\rho'$.
\end{theorem}

\begin{theorem}[Restriction assignment weight at bad extension le]
\lean{FABL.DNFFormula.restrictionAssignmentWeightAt_badExtension_le}
\statementsource{arora-barak-ch13-circuit-lowerbounds}{pdf-52929dbaa962-p003-b004}
\proofsource{Analysis-of-Boolean-Functions-by-Ryan-ODonnell}{pdf-ffbe9c7b3306-p110-b003}
\label{FABL.DNFFormula.restrictionAssignmentWeightAt_badExtension_le}
\leanok
\uses{FABL.CoordRestriction, FABL.DNFFormula, FABL.DNFFormula.IsBadRestriction, FABL.DNFFormula.badExtension, FABL.DNFFormula.restrictionAssignmentWeightAt_badExtension_eq, FABL.restrictionAssignmentWeightAt, FABL.restrictionAssignmentWeightAt_nonneg}
\difficulty{4}
% dataset-pipeline: stub -- prose pending build_informal_statements.py
Statement pending regeneration.
\end{theorem}

\begin{definition}[Bad restriction weight]
\lean{FABL.DNFFormula.badRestrictionWeight}
\statementsource{arora-barak-ch13-circuit-lowerbounds}{pdf-52929dbaa962-p003-b004}
\statementsource{Analysis-of-Boolean-Functions-by-Ryan-ODonnell}{pdf-ffbe9c7b3306-p100-b004}
\label{FABL.DNFFormula.badRestrictionWeight}
\leanok
\uses{FABL.CoordRestriction, FABL.DNFFormula, FABL.DNFFormula.IsBadRestriction, FABL.restrictionAssignmentWeightAt}
Let $\varphi$ be a DNF formula on $n$ coordinates and let $\delta\in\mathbb R$. Its bad
restriction weight is the sum, over all restrictions on which $\varphi$ is nonconstant,
of the assignment weight at parameter $\delta$ of that restriction.
\end{definition}

\begin{theorem}[Sum weight bad extension le]
\lean{FABL.DNFFormula.sum_weight_badExtension_le}
\statementsource{arora-barak-ch13-circuit-lowerbounds}{
  pdf-52929dbaa962-p003-b004,
  pdf-52929dbaa962-p004-b002,
  pdf-52929dbaa962-p005-b001
}
\proofsource{Analysis-of-Boolean-Functions-by-Ryan-ODonnell}{pdf-ffbe9c7b3306-p110-b003}
\label{FABL.DNFFormula.sum_weight_badExtension_le}
\leanok
\uses{FABL.CoordRestriction, FABL.DNFFormula, FABL.DNFFormula.IsBadRestriction, FABL.DNFFormula.badExtensionFiber, FABL.DNFFormula.badExtensionFiber_card_le_width, FABL.DNFFormula.badExtensionMap, FABL.DNFFormula.width, FABL.restrictionAssignmentWeightAt, FABL.restrictionAssignmentWeightAt_nonneg, FABL.sum_restrictionAssignmentWeightAt}
\difficulty{4}
% dataset-pipeline: stub -- prose pending build_informal_statements.py
Statement pending regeneration.
\end{theorem}

\begin{theorem}[Bad restriction weight for bounded-width DNF]
\lean{FABL.exercise4_19}
\proofsource{Analysis-of-Boolean-Functions-by-Ryan-ODonnell}{pdf-ffbe9c7b3306-p110-b003}
\label{FABL.exercise4_19}
\leanok
\uses{FABL.CoordRestriction, FABL.DNFFormula, FABL.DNFFormula.IsBadRestriction, FABL.DNFFormula.badExtensionMap, FABL.DNFFormula.badRestrictionWeight, FABL.DNFFormula.restrictionAssignmentWeightAt_badExtension_le, FABL.DNFFormula.sum_weight_badExtension_le, FABL.DNFFormula.width, FABL.restrictionAssignmentWeightAt}
\difficulty{4}
Let $\varphi$ be a disjunctive normal form on $n$ Boolean variables whose terms each
contain distinct variables, and suppose its width is at most $w$. Choose each coordinate
restriction independently: leave it free with weight $\delta$, or fix it to either sign
with weight $(1-\delta)/2$. Call a restriction bad when the restricted function is
nonconstant. If $0\le\delta\le 1/3$, then the total weight of bad restrictions is at
most $3\delta w$.
\end{theorem}

\begin{theorem}[Evaluation of a depth-zero tree]
\lean{FABL.F₂DecisionTree.eval_eq_of_depth_eq_zero}
\leanok
\uses{FABL.DecisionTree, FABL.F₂DecisionTree.depth, FABL.F₂DecisionTree.eval, FABL.Sign}
\difficulty{4}
Let $n$ be a nonnegative integer and let $T$ be a decision tree on the $n$-dimensional
Boolean cube whose outputs are signs in $\{ -1,1\}$. If the depth of $T$ is zero, then
$T$ has the same evaluation on every two inputs $x,y\in\{0,1\}^n$.
\end{theorem}

\begin{theorem}[Decision tree depth detects nonconstancy]
\lean{FABL.one_le_decisionTreeDepth_iff_nonconstant}
\label{FABL.one_le_decisionTreeDepth_iff_nonconstant}
\leanok
\uses{FABL.DecisionTree, FABL.F₂DecisionTree.Computes, FABL.F₂DecisionTree.computes_iff, FABL.F₂DecisionTree.decisionTreeDepth, FABL.F₂DecisionTree.decisionTreeDepth_le_of_computes, FABL.F₂DecisionTree.depth, FABL.F₂DecisionTree.eval, FABL.F₂DecisionTree.eval_eq_of_depth_eq_zero, FABL.F₂DecisionTree.exists_computingTree_depth_eq_decisionTreeDepth, FABL.Sign}
\difficulty{5}
Let $f:\{0,1\}^n\to\{−1,1\}$ be a Boolean function on the $n$-dimensional Boolean cube.
Its minimum decision tree depth is at least $1$ if and only if there exist inputs $x$
and $y$ such that $f(x)\ne f(y)$.
\end{theorem}

\begin{theorem}[Failure indicator and bad restrictions]
\lean{FABL.coordSwitchingFailureIndicator_one_eq_badRestriction}
\label{FABL.coordSwitchingFailureIndicator_one_eq_badRestriction}
\leanok
\uses{FABL.CoordRestriction, FABL.CoordRestriction.restrict, FABL.DNFFormula, FABL.DNFFormula.IsBadRestriction, FABL.DNFFormula.toBooleanFunction, FABL.coordSwitchingFailureIndicator, FABL.one_le_decisionTreeDepth_iff_nonconstant}
\difficulty{4}
Let $\phi$ be a disjunctive normal form on the sign cube $\{-1,1\}^n$, and let
$\rho:[n]\to\{\text{free},\text{fixed }1,\text{fixed }-1\}$ be a coordinate
restriction. The coordinate-switching failure indicator at threshold $1$ for the
restriction of $\phi$ equals $1$ if there are two Boolean inputs on which the
restricted function takes different values, and equals $0$ otherwise.
\end{theorem}

\begin{theorem}[Bad restriction weight and switching probability]
\lean{FABL.DNFFormula.badRestrictionWeight_eq_coordSwitchingFailureProbability}
\label{FABL.DNFFormula.badRestrictionWeight_eq_coordSwitchingFailureProbability}
\leanok
\uses{FABL.DNFFormula, FABL.DNFFormula.badRestrictionWeight, FABL.DNFFormula.toBooleanFunction, FABL.coordSwitchingFailureIndicator_one_eq_badRestriction, FABL.coordSwitchingFailureProbability}
\difficulty{3}
Let $n$ be a natural number, let $\varphi$ be a disjunctive normal form on $n$ sign
variables, and let $\delta$ be real. The total weight, under the product restriction
weights with free-coordinate weight $\delta$ and each fixed-sign weight $(1-\delta)/2$,
of restrictions on which $\varphi$ is nonconstant equals the coordinate-switching
failure probability at threshold $1$ for the Boolean function represented by $\varphi$.
\end{theorem}

\begin{theorem}[Switching failure probability for bounded width]
\lean{FABL.exercise4_19_switchingFailureProbability}
\proofsource{Analysis-of-Boolean-Functions-by-Ryan-ODonnell}{pdf-ffbe9c7b3306-p110-b003}
\label{FABL.exercise4_19_switchingFailureProbability}
\leanok
\uses{FABL.DNFFormula, FABL.DNFFormula.badRestrictionWeight_eq_coordSwitchingFailureProbability, FABL.DNFFormula.toBooleanFunction, FABL.DNFFormula.width, FABL.coordSwitchingFailureProbability_eq, FABL.exercise4_19, FABL.switchingFailureProbability}
\difficulty{3}
Let $\varphi$ be a disjunctive normal form on the $n$-dimensional Boolean cube whose
terms each contain at most $w$ literals. For $0\leq\delta\leq 1/3$, choose each
coordinate independently to be fixed with probability $1-\delta$ and otherwise left
free, and choose the fixed signs uniformly; then the probability that the restricted
function has decision-tree depth at least $1$ is at most $3\delta w$.
\end{theorem}

\begin{theorem}[Constancy of width-zero formulas]
\lean{FABL.DNFFormula.toBooleanFunction_eq_of_width_eq_zero}
\label{FABL.DNFFormula.toBooleanFunction_eq_of_width_eq_zero}
\leanok
\uses{FABL.DNFFormula, FABL.DNFFormula.eval, FABL.DNFFormula.eval_eq_neg_one_iff, FABL.DNFFormula.toBooleanFunction, FABL.DNFFormula.width, FABL.DNFFormula.width_le_of_mem, FABL.DNFTerm.eval, FABL.DNFTerm.eval_eq_neg_one_iff, FABL.DNFTerm.width}
\difficulty{4}
Let $n\in\mathbb{N}$, and let $\varphi$ be a DNF formula on $n$ coordinates with width
zero. Its associated Boolean function takes the same value on every pair of inputs in
the sign cube.
\end{theorem}

\begin{theorem}[Zero bad restriction weight at width zero]
\lean{FABL.DNFFormula.badRestrictionWeight_eq_zero_of_width_eq_zero}
\label{FABL.DNFFormula.badRestrictionWeight_eq_zero_of_width_eq_zero}
\leanok
\uses{FABL.CoordRestriction.complete, FABL.DNFFormula, FABL.DNFFormula.badRestrictionWeight, FABL.DNFFormula.toBooleanFunction_eq_of_width_eq_zero, FABL.DNFFormula.width}
\difficulty{4}
Let $n$ be a natural number, let $\varphi$ be a disjunctive normal form on $n$ sign
variables, and suppose its width is zero. For every real $\delta$, the total weight of
restrictions on which $\varphi$ is nonconstant is zero.
\end{theorem}

\begin{theorem}[Baby switching lemma dnf]
\lean{FABL.babySwitchingLemma_dnf}
\statementsource{lmn-constant-depth-fourier}{pdf-63759bd909d5-p005-b004}
\statementsource{arora-barak-ch13-circuit-lowerbounds}{
  pdf-52929dbaa962-p002-b003,
  pdf-52929dbaa962-p003-b004,
  pdf-52929dbaa962-p003-b005,
  pdf-52929dbaa962-p004-b001,
  pdf-52929dbaa962-p004-b002,
  pdf-52929dbaa962-p005-b001
}
\proofsource{Analysis-of-Boolean-Functions-by-Ryan-ODonnell}{
  pdf-ffbe9c7b3306-p100-b004,
  pdf-ffbe9c7b3306-p110-b003
}
\label{FABL.babySwitchingLemma_dnf}
\leanok
\uses{FABL.DNFFormula, FABL.DNFFormula.badRestrictionWeight_eq_coordSwitchingFailureProbability, FABL.DNFFormula.badRestrictionWeight_eq_zero_of_width_eq_zero, FABL.DNFFormula.toBooleanFunction, FABL.DNFFormula.width, FABL.coordSwitchingFailureProbability_eq, FABL.exercise4_19_switchingFailureProbability, FABL.switchingFailureProbability, FABL.switchingFailureProbability_le_one}
\difficulty{5}
% dataset-pipeline: stub -- prose pending build_informal_statements.py
Statement pending regeneration.
\end{theorem}

\begin{theorem}[Addition in the two-element field]
\lean{FABL.f₂_one_add_one}
\leanok
\uses{FABL.𝔽₂}
\difficulty{1}
In the field $br/2br$, the sum $1+1$ is $0$.
\end{theorem}

\begin{theorem}[Involution of the binary-cube complement]
\lean{FABL.binaryCubeComplement_involutive}
\label{FABL.binaryCubeComplement_involutive}
\leanok
\difficulty{3}
For every natural number $n$ and every binary vector $x$ of length $n$, applying the
binary-cube complement twice returns $x$.
\end{theorem}

\begin{theorem}[Sign encoding under coordinate complementation]
\lean{FABL.signEncode_binaryCubeComplement}
\label{FABL.signEncode_binaryCubeComplement}
\leanok
\uses{FABL.signEncode, FABL.signEncode_add}
\difficulty{3}
Let $n$ be a natural number, let $x$ be a vector of length $n$ over the two-element
field, and let $i$ be one of its coordinates. The sign encoding of the vector obtained
by complementing coordinate $i$ is the negative of the sign encoding of the original
$i$th coordinate.
\end{theorem}

\begin{definition}[Complementing decision tree inputs]
\lean{FABL.F₂DecisionTree.complementInputs}
\leanok
\uses{FABL.F₂DecisionTree}
Let a decision tree on $n$ binary coordinates, with values in a type $\alpha$, query
only coordinates in an available set. Complementing its inputs leaves each leaf value
unchanged and, at every query, exchanges the two child subtrees while retaining the
queried coordinate and the available set.
\end{definition}

\begin{theorem}[Evaluation under complemented inputs]
\lean{FABL.F₂DecisionTree.eval_complementInputs}
\leanok
\uses{FABL.F₂DecisionTree, FABL.F₂DecisionTree.complementInputs, FABL.F₂DecisionTree.eval}
\difficulty{4}
Let $n$ be a nonnegative integer, let $I$ be a set of coordinates among
$1,\ldots,n$, let $\alpha$ be a type, and let $T$ be a binary decision tree with
output values in $\alpha$ that may query only coordinates in $I$. For every Boolean
input $x\in\{0,1\}^n$, evaluating the tree with its zero and one branches interchanged
at each query gives the same result as evaluating the original tree on the
coordinatewise complement of $x$.
\end{theorem}

\begin{theorem}[Depth under input complementation]
\lean{FABL.F₂DecisionTree.depth_complementInputs}
\leanok
\uses{FABL.F₂DecisionTree, FABL.F₂DecisionTree.complementInputs, FABL.F₂DecisionTree.depth}
\difficulty{4}
For any set of available coordinates and any binary decision tree $T$ with those
coordinates available, complementing the input bits leaves its depth unchanged: the
complemented tree has depth equal to the depth of $T$.
\end{theorem}

\begin{definition}[Mapping decision tree outputs]
\lean{FABL.F₂DecisionTree.mapOutputs}
\leanok
\uses{FABL.F₂DecisionTree}
Let $g$ be a function from a type $\alpha$ to a type $\beta$. Applying $g$ to the
outputs of a decision tree replaces each leaf value $a$ by $g(a)$ and preserves every
query and its two child subtrees, producing a tree with leaf values in $\beta$.
\end{definition}

\begin{theorem}[Evaluation after mapping outputs]
\lean{FABL.F₂DecisionTree.eval_mapOutputs}
\leanok
\uses{FABL.F₂DecisionTree, FABL.F₂DecisionTree.eval, FABL.F₂DecisionTree.mapOutputs}
\difficulty{4}
Let $n$ be a nonnegative integer, let $alpha$ and $beta$ be types, let
$g:alpha\tobeta$, and let $T$ be a binary decision tree on $n$ coordinates with
outputs in $alpha$. For every input $x\in(\mathbb{F}_2)^n$, evaluating the tree
obtained by applying $g$ to each leaf output gives $g$ applied to the evaluation of $T$
at $x$.
\end{theorem}

\begin{theorem}[Depth under output mapping]
\lean{FABL.F₂DecisionTree.depth_mapOutputs}
\leanok
\uses{FABL.F₂DecisionTree, FABL.F₂DecisionTree.depth, FABL.F₂DecisionTree.mapOutputs}
\difficulty{4}
Let $g$ be any map from one output type to another. For any binary decision tree with a
specified finite set of available coordinates, applying $g$ to every leaf output leaves
the tree's depth unchanged.
\end{theorem}

\begin{definition}[Binary Boolean dual]
\lean{FABL.binaryBooleanDual}
\statementsource{Analysis-of-Boolean-Functions-by-Ryan-ODonnell}{pdf-ffbe9c7b3306-p015-b006}
\label{FABL.binaryBooleanDual}
\leanok
\uses{FABL.Sign}
For a function $f:\mathbb F_2^n\to\{\pm1\}$, define its binary Boolean dual by sending
each $x\in\mathbb F_2^n$ to the negative of $f$ evaluated at the coordinatewise
complement of $x$.
\end{definition}

\begin{theorem}[One-sided depth bound under Boolean duality]
\lean{FABL.decisionTreeDepth_binaryBooleanDual_le}
\label{FABL.decisionTreeDepth_binaryBooleanDual_le}
\leanok
\uses{FABL.DecisionTree, FABL.F₂DecisionTree.Computes, FABL.F₂DecisionTree.complementInputs, FABL.F₂DecisionTree.computes_iff, FABL.F₂DecisionTree.decisionTreeDepth, FABL.F₂DecisionTree.decisionTreeDepth_le_of_computes, FABL.F₂DecisionTree.depth, FABL.F₂DecisionTree.eval_complementInputs, FABL.F₂DecisionTree.eval_mapOutputs, FABL.F₂DecisionTree.exists_computingTree_depth_eq_decisionTreeDepth, FABL.F₂DecisionTree.mapOutputs, FABL.Sign, FABL.binaryBooleanDual}
\difficulty{4}
Let $n\in\mathbb{N}$ and let $f:\mathbb{F}_2^n\to\{−1,1\}$. Define
$f^*(x)=-f(\mathbf{1}-x)$, with subtraction taken coordinatewise in $\mathbb{F}_2^n$.
The minimum depth of a binary decision tree computing $f^*$ is at most the minimum depth
of one computing $f$.
\end{theorem}

\begin{theorem}[Involutivity of Boolean duality]
\lean{FABL.binaryBooleanDual_involutive}
\label{FABL.binaryBooleanDual_involutive}
\leanok
\uses{FABL.Sign, FABL.binaryBooleanDual}
\difficulty{3}
Let $n$ be a natural number and let $f$ map binary vectors of length $n$ to the signs
$71$. Define its dual by negating the value of $f$ at the coordinatewise complement of
the input. Taking the dual twice returns $f$.
\end{theorem}

\begin{theorem}[Duality preserves decision tree depth]
\lean{FABL.decisionTreeDepth_binaryBooleanDual}
\label{FABL.decisionTreeDepth_binaryBooleanDual}
\leanok
\uses{FABL.F₂DecisionTree.decisionTreeDepth, FABL.Sign, FABL.binaryBooleanDual, FABL.binaryBooleanDual_involutive, FABL.decisionTreeDepth_binaryBooleanDual_le}
\difficulty{4}
Let $n\in\mathbb{N}$ and let $f:\mathbb{F}_2^n\to\{−1,1\}$. Define its dual by
$f^*(x)=-f(\mathbf{1}-x)$, with subtraction taken coordinatewise in $\mathbb{F}_2^n$.
Then the minimum depth of a binary decision tree computing $f^*$ equals the minimum
depth of a binary decision tree computing $f$.
\end{theorem}

\begin{definition}[Negating a restriction state]
\lean{FABL.CoordRestriction.negateState}
\label{FABL.CoordRestriction.negateState}
\leanok
\uses{FABL.CoordRestriction}
Negation leaves a free coordinate free, changes a coordinate fixed at $1$ to one fixed
at $-1$, and changes a coordinate fixed at $-1$ to one fixed at $1$.
\end{definition}

\begin{theorem}[Negating a restriction state twice]
\lean{FABL.CoordRestriction.negateState_involutive}
\label{FABL.CoordRestriction.negateState_involutive}
\leanok
\uses{FABL.CoordRestriction, FABL.CoordRestriction.negateState}
\difficulty{4}
Negating any restriction state twice returns that state.
\end{theorem}

\begin{definition}[Negating a restriction assignment]
\lean{FABL.CoordRestriction.negateAssignment}
\label{FABL.CoordRestriction.negateAssignment}
\leanok
\uses{FABL.CoordRestriction, FABL.CoordRestriction.negateState}
Given an assignment of restriction states to the $n$ coordinates, its negation is the
assignment obtained by negating the state at every coordinate: free states remain free,
and the two fixed states are exchanged.
\end{definition}

\begin{theorem}[Negating an assignment twice]
\lean{FABL.CoordRestriction.negateAssignment_involutive}
\label{FABL.CoordRestriction.negateAssignment_involutive}
\leanok
\uses{FABL.CoordRestriction, FABL.CoordRestriction.negateAssignment}
\difficulty{3}
Let $n$ be a natural number and let $\rho$ assign a restriction state to each of the $n$
coordinates. Negating every state in $\rho$ twice returns $\rho$.
\end{theorem}

\begin{definition}[Negation of restriction assignments]
\lean{FABL.coordRestrictionNegationEquiv}
\label{FABL.coordRestrictionNegationEquiv}
\leanok
\uses{FABL.CoordRestriction, FABL.CoordRestriction.negateAssignment, FABL.CoordRestriction.negateAssignment_involutive}
For each $n$, negating every state defines a bijection from the set of restriction
assignments on the $n$ coordinates to itself. This bijection is its own inverse.
\end{definition}

\begin{theorem}[Invariance of coordinate weights under state negation]
\lean{FABL.coordRestrictionWeightAt_negateState}
\label{FABL.coordRestrictionWeightAt_negateState}
\leanok
\uses{FABL.CoordRestriction, FABL.CoordRestriction.negateState, FABL.coordRestrictionWeightAt}
\difficulty{4}
For every real number $\delta$ and every coordinate restriction $c$, the weight assigned
to the restriction obtained by negating its fixed state equals the weight assigned to
$c$.
\end{theorem}

\begin{theorem}[Negation invariance of assignment weight]
\lean{FABL.restrictionAssignmentWeightAt_negateAssignment}
\label{FABL.restrictionAssignmentWeightAt_negateAssignment}
\leanok
\uses{FABL.CoordRestriction, FABL.CoordRestriction.negateAssignment, FABL.restrictionAssignmentWeightAt}
\difficulty{1}
Let $n$ be a natural number, let $
ho$ assign to each of the $n$ coordinates one of the
states “free,” “fixed to $1$,” or “fixed to $-1$,” and let $
ho^{-}$ be obtained by
exchanging the two fixed states while leaving free states unchanged. For every real
parameter $b4$, the product weight of $
ho^{-}$, using coordinate weights $b4$ for a
free state and $(1-b4)/2$ for either fixed state, equals the product weight of $
ho$.
\end{theorem}

\begin{theorem}[Completion under negated restrictions]
\lean{FABL.CoordRestriction.complete_negateAssignment}
\label{FABL.CoordRestriction.complete_negateAssignment}
\leanok
\uses{FABL.CoordRestriction, FABL.CoordRestriction.complete, FABL.CoordRestriction.negateAssignment, FABL.CoordRestriction.negateState, FABL.signEncode_binaryCubeComplement}
\difficulty{4}
Let $n$ be a natural number, let $\rho$ assign a restriction state to each of the $n$
coordinates, and let $x\in\mathbb{F}_2^n$. Completing the bitwise complement of $x$
under the restriction assignment obtained by negating every state in $\rho$ gives the
coordinatewise negative of the completion of $x$ under $\rho$.
\end{theorem}

\begin{theorem}[Binary duality and restriction]
\lean{FABL.CoordRestriction.restrict_booleanDual}
\label{FABL.CoordRestriction.restrict_booleanDual}
\leanok
\uses{FABL.BooleanFunction, FABL.CNFFormula.booleanDual, FABL.CoordRestriction, FABL.CoordRestriction.complete_negateAssignment, FABL.CoordRestriction.negateAssignment, FABL.CoordRestriction.restrict, FABL.binaryBooleanDual}
\difficulty{3}
Let $n$ be a nonnegative integer, let $f$ be a function from the $n$-dimensional sign
cube to $\{\pm1\}$, and let $\rho$ assign each coordinate one of the states free, fixed
to $1$, or fixed to $-1$. Restricting the Boolean dual of $f$ according to $\rho$ gives
the binary Boolean dual of the restriction of $f$ according to the coordinatewise
negation of $\rho$.
\end{theorem}

\begin{theorem}[Depth under duality and coordinate restrictions]
\lean{FABL.decisionTreeDepth_restrict_booleanDual}
\label{FABL.decisionTreeDepth_restrict_booleanDual}
\leanok
\uses{FABL.BooleanFunction, FABL.CNFFormula.booleanDual, FABL.CoordRestriction, FABL.CoordRestriction.negateAssignment, FABL.CoordRestriction.restrict, FABL.CoordRestriction.restrict_booleanDual, FABL.F₂DecisionTree.decisionTreeDepth, FABL.decisionTreeDepth_binaryBooleanDual}
\difficulty{1}
Let $n\in\mathbb{N}$, let $f:\{−1,1\}^n\to\{−1,1\}$, and assign to each coordinate one
of three states: free, fixed to $1$, or fixed to $-1$. For each binary input, form a
sign vector by encoding each free bit as $1$ or $-1$ and using the assigned signs at
fixed coordinates; restrict $f$ by evaluating it on this vector. If the dual function is
$f^*(x)=-f(-x)$, then for every assignment of coordinate states, the restricted dual has
the same minimum decision tree depth as the restriction of $f$ obtained by interchanging
the two fixed states and leaving free states unchanged.
\end{theorem}

\begin{theorem}[Switching failure under Boolean duality]
\lean{FABL.coordSwitchingFailureIndicator_booleanDual}
\label{FABL.coordSwitchingFailureIndicator_booleanDual}
\leanok
\uses{FABL.BooleanFunction, FABL.CNFFormula.booleanDual, FABL.CoordRestriction, FABL.CoordRestriction.negateAssignment, FABL.coordSwitchingFailureIndicator, FABL.decisionTreeDepth_restrict_booleanDual}
\difficulty{1}
Let $f:\{\pm1\}^n\to\{\pm1\}$ be a Boolean function, let $k$ be a natural number, and
let $\rho$ assign a coordinate restriction to each of the $n$ coordinates. Define the
Boolean dual by $f^{*}(x)=-f(-x)$, and let $-\rho$ denote the assignment obtained by
negating each fixed state while leaving free coordinates free. The switching failure
indicator of $f^{*}$ at threshold $k$ and assignment $\rho$ equals that of $f$ at
threshold $k$ and assignment $-\rho$.
\end{theorem}

\begin{theorem}[Boolean duality preserves failure probability]
\lean{FABL.coordSwitchingFailureProbability_booleanDual}
\label{FABL.coordSwitchingFailureProbability_booleanDual}
\leanok
\uses{FABL.BooleanFunction, FABL.CNFFormula.booleanDual, FABL.CoordRestriction, FABL.CoordRestriction.negateAssignment, FABL.coordRestrictionNegationEquiv, FABL.coordSwitchingFailureIndicator, FABL.coordSwitchingFailureIndicator_booleanDual, FABL.coordSwitchingFailureProbability, FABL.restrictionAssignmentWeightAt, FABL.restrictionAssignmentWeightAt_negateAssignment}
\difficulty{4}
Let $f:\{-1,1\}^n\to\{\pm1\}$ be a Boolean function, let $\delta\in\mathbb R$,
and let $k$ be a natural number. The coordinate-switching failure probability at
parameter $\delta$ and threshold $k$ is unchanged when $f$ is replaced by its Boolean
dual, the function $x\mapsto -f(-x)$.
\end{theorem}

\begin{theorem}[Dual invariance of switching failure]
\lean{FABL.switchingFailureProbability_booleanDual}
\label{FABL.switchingFailureProbability_booleanDual}
\leanok
\uses{FABL.BooleanFunction, FABL.CNFFormula.booleanDual, FABL.coordSwitchingFailureProbability_booleanDual, FABL.coordSwitchingFailureProbability_eq, FABL.switchingFailureProbability}
\difficulty{3}
Let $f:\{-1,1\}^n\to\{-1,1\}$ be a sign-valued Boolean function, let
$\delta\in\mathbb R$, and let $k\in\mathbb N$. Define its Boolean dual by
$f^*(x)=-f(-x)$. The switching-failure probability at parameters $\delta,k$ is
unchanged by dualization: $P_{\delta,k}(f^*)=P_{\delta,k}(f)$, where
$P_{\delta,k}(f)$ is the weighted average over random restrictions of the indicator
that the restricted function has decision-tree depth at least $k$.
\end{theorem}

\begin{theorem}[Baby switching lemma]
\lean{FABL.babySwitchingLemma}
\statementsource{lmn-constant-depth-fourier}{pdf-63759bd909d5-p005-b004}
\statementsource{arora-barak-ch13-circuit-lowerbounds}{
  pdf-52929dbaa962-p002-b003,
  pdf-52929dbaa962-p003-b004,
  pdf-52929dbaa962-p003-b005,
  pdf-52929dbaa962-p004-b002,
  pdf-52929dbaa962-p005-b001
}
\proofsource{Analysis-of-Boolean-Functions-by-Ryan-ODonnell}{pdf-ffbe9c7b3306-p100-b004}
\label{FABL.babySwitchingLemma}
\leanok
\uses{FABL.BooleanFunction, FABL.CNFFormula.booleanDual, FABL.CNFFormula.switchAndOr, FABL.CNFFormula.switchAndOr_toBooleanFunction, FABL.CNFFormula.toBooleanFunction, FABL.CNFFormula.width, FABL.DNFFormula.toBooleanFunction, FABL.DNFFormula.width, FABL.HasCNFWidthLE, FABL.HasDNFWidthLE, FABL.babySwitchingLemma_dnf, FABL.switchingFailureProbability, FABL.switchingFailureProbability_booleanDual}
\difficulty{5}
% dataset-pipeline: stub -- prose pending build_informal_statements.py
Statement pending regeneration.
\end{theorem}

\begin{definition}[Embedding a decision tree]
\lean{FABL.F₂DecisionTree.pushEmbedding}
\leanok
\uses{FABL.F₂DecisionTree}
Let $m,n$ be natural numbers, let $\alpha$ be a type, and let $e:[m]\hookrightarrow[n]$
be an embedding. For finite coordinate sets $S\subseteq[m]$ and $T\subseteq[n]$ such
that $e(i)\in T$ for every $i\in S$, every binary decision tree on $[m]$ with available
coordinates $S$ induces a binary decision tree on $[n]$ with available coordinates $T$:
leaves retain their values, and each query at $i$ becomes a query at $e(i)$, with the
two branches embedded recursively.
\end{definition}

\begin{theorem}[Depth under pushing forward by an embedding]
\lean{FABL.F₂DecisionTree.depth_pushEmbedding}
\leanok
\uses{FABL.F₂DecisionTree, FABL.F₂DecisionTree.depth, FABL.F₂DecisionTree.pushEmbedding}
\difficulty{4}
Let $e$ be an embedding of the coordinates of a finite cube of dimension $m$ into one of
dimension $n$. Let $S$ and $T$ be sets of available coordinates in the two cubes such
that $e(i)\in T$ for every $i\in S$. For any binary decision tree on $S$, pushing it
forward along $e$ preserves its depth.
\end{theorem}

\begin{theorem}[Evaluation under coordinate embedding]
\lean{FABL.F₂DecisionTree.eval_pushEmbedding}
\leanok
\uses{FABL.F₂DecisionTree, FABL.F₂DecisionTree.eval, FABL.F₂DecisionTree.pushEmbedding}
\difficulty{4}
Let $m,n\in\mathbb{N}$, let $e:\{0,\ldots,m-1\}\hookrightarrow\{0,\ldots,n-1\}$ be an
injection, and let $S\subseteq\{0,\ldots,m-1\}$ and $U\subseteq\{0,\ldots,n-1\}$ satisfy
$e(i)\in U$ for every $i\in S$. For a decision tree $T$ on the $m$-coordinate cube with
available set $S$, and every input $x\in(\mathbb{F}_2)^n$, evaluating the tree embedded
along $e$ at $x$ equals evaluating $T$ on the input whose $i$th coordinate is
$x_{e(i)}$.
\end{theorem}

\begin{definition}[Pulling back a decision tree]
\lean{FABL.F₂DecisionTree.pullEmbedding}
\leanok
\uses{FABL.F₂DecisionTree}
Let $m,n$ be natural numbers, let $\alpha$ be a type, and let $e:[m]\hookrightarrow[n]$
be an embedding. For finite coordinate sets $S\subseteq[m]$ and $T\subseteq[n]$ such
that $i\in S$ exactly when $e(i)\in T$, every binary decision tree on $[n]$ with
available coordinates $T$ induces a binary decision tree on $[m]$ with available
coordinates $S$: a query at a coordinate in the image of $e$ becomes the query at its
unique preimage, while a query outside that image is discarded and the zero branch is
retained; leaves retain their values.
\end{definition}

\begin{theorem}[Depth under pullback by an embedding]
\lean{FABL.F₂DecisionTree.depth_pullEmbedding_le}
\leanok
\uses{FABL.F₂DecisionTree, FABL.F₂DecisionTree.depth, FABL.F₂DecisionTree.pullEmbedding}
\difficulty{5}
Let $e$ be an embedding of the coordinates of a finite cube of dimension $m$ into one of
dimension $n$. Let $S$ and $T$ be sets of available coordinates in the two cubes such
that, for every coordinate $i$ of the first cube, $i$ belongs to $S$ exactly when $e(i)$
belongs to $T$. For any binary decision tree on $T$, its pullback along $e$ has depth at
most the original tree's depth.
\end{theorem}

\begin{theorem}[Evaluation under coordinate pullback]
\lean{FABL.F₂DecisionTree.eval_pullEmbedding}
\leanok
\uses{FABL.F₂DecisionTree, FABL.F₂DecisionTree.eval, FABL.F₂DecisionTree.pullEmbedding}
\difficulty{6}
Let $m,n\in\mathbb{N}$, let $e:\{0,\ldots,m-1\}\hookrightarrow\{0,\ldots,n-1\}$ be an
injection, and let $S\subseteq\{0,\ldots,m-1\}$ and $U\subseteq\{0,\ldots,n-1\}$ satisfy
$i\in S$ if and only if $e(i)\in U$ for every $i$. For a decision tree $T$ on the
$n$-coordinate cube with available set $U$, and every input $x\in(\mathbb{F}_2)^m$,
evaluating the tree pulled back along $e$ at $x$ equals evaluating $T$ on the
$n$-coordinate input that takes value $x_i$ at $e(i)$ and $0$ at every coordinate
outside the image of $e$.
\end{theorem}

\begin{theorem}[Decision tree depth under coordinate embedding]
\lean{FABL.F₂DecisionTree.decisionTreeDepth_comp_embedding}
\leanok
\uses{FABL.DecisionTree, FABL.F₂DecisionTree.Computes, FABL.F₂DecisionTree.computes_iff, FABL.F₂DecisionTree.decisionTreeDepth, FABL.F₂DecisionTree.decisionTreeDepth_le_of_computes, FABL.F₂DecisionTree.depth, FABL.F₂DecisionTree.depth_pullEmbedding_le, FABL.F₂DecisionTree.depth_pushEmbedding, FABL.F₂DecisionTree.eval, FABL.F₂DecisionTree.eval_pullEmbedding, FABL.F₂DecisionTree.eval_pushEmbedding, FABL.F₂DecisionTree.exists_computingTree_depth_eq_decisionTreeDepth, FABL.F₂DecisionTree.pullEmbedding, FABL.F₂DecisionTree.pushEmbedding}
\difficulty{6}
Let $m,n$ be nonnegative integers, let
$e:\{0,\ldots,m-1\}\hookrightarrow\{0,\ldots,n-1\}$ be an injection, and let $f:\mathbb
F_2^m\to\alpha$. Define a function on $\mathbb F_2^n$ by reading the coordinates indexed
by $e$ and applying $f$. Its decision tree depth equals the decision tree depth of $f$.
\end{theorem}

\begin{theorem}[Decision tree depth under coordinate equivalence]
\lean{FABL.F₂DecisionTree.decisionTreeDepth_comp_equiv}
\leanok
\uses{FABL.F₂DecisionTree.decisionTreeDepth, FABL.F₂DecisionTree.decisionTreeDepth_comp_embedding}
\difficulty{1}
Let $m,n$ be nonnegative integers, let $e:\{0,\ldots,m-1\}\cong\{0,\ldots,n-1\}$ be a
bijection, and let $f:\mathbb F_2^m\to\alpha$. Define a function on $\mathbb F_2^n$ by
reading the coordinates indexed by $e$ and applying $f$. Its decision tree depth equals
the decision tree depth of $f$.
\end{theorem}

\begin{definition}[Reindexing a Boolean function]
\lean{FABL.reindexBooleanFunction}
\label{FABL.reindexBooleanFunction}
\leanok
\uses{FABL.BooleanFunction}
Let $m,n$ be natural numbers and let $e:[m]\simeq[n]$ be a bijection. Given a function
$f$ from the $m$-dimensional Boolean cube to the sign group, define its reindexing on
the $n$-dimensional Boolean cube by assigning to each input $x$ the value of $f$ at the
input whose $i$th coordinate is $x(e(i))$.
\end{definition}

\begin{theorem}[Evaluation After Boolean-Function Reindexing]
\lean{FABL.reindexBooleanFunction_apply}
\label{FABL.reindexBooleanFunction_apply}
\leanok
\uses{FABL.BooleanFunction, FABL.reindexBooleanFunction}
\difficulty{1}
Let $m,n$ be natural numbers, let $e\colon\{0,\ldots,m-1\}\simeq\{0,\ldots,n-1\}$ be a
bijection, let $f\colon\{-1,1\}^m\to\{-1,1\}$, and let $x\in\{-1,1\}^n$. The reindexed
function evaluated at $x$ equals $f$ evaluated on the input whose $i$th coordinate is
$x_{e(i)}$.
\end{theorem}

\begin{theorem}[Restriction under coordinate reindexing]
\lean{FABL.CoordRestriction.restrict_reindexBooleanFunction}
\label{FABL.CoordRestriction.restrict_reindexBooleanFunction}
\leanok
\uses{FABL.BooleanFunction, FABL.CoordRestriction, FABL.CoordRestriction.restrict, FABL.reindexBooleanFunction}
\difficulty{1}
Let $m,n$ be nonnegative integers, let $e$ be a bijection from the coordinates in $[m]$
to those in $[n]$, let $f$ be a function from the $m$-dimensional sign cube to
$\{\pm1\}$, and let $\rho$ assign a restriction state to each coordinate in $[n]$. For
every binary input $x$ on $[n]$, restricting the reindexed function according to $\rho$
equals evaluating the restriction of $f$ whose coordinate states are $\rho(e(i))$ at the
input whose $i$th coordinate is $x(e(i))$.
\end{theorem}

\begin{theorem}[Reindexing restricted decision tree depth]
\lean{FABL.restrictedDecisionTreeDepth_reindexBooleanFunction}
\label{FABL.restrictedDecisionTreeDepth_reindexBooleanFunction}
\leanok
\uses{FABL.BooleanFunction, FABL.CoordRestriction, FABL.CoordRestriction.restrict, FABL.CoordRestriction.restrict_reindexBooleanFunction, FABL.F₂DecisionTree.decisionTreeDepth, FABL.F₂DecisionTree.decisionTreeDepth_comp_equiv, FABL.reindexBooleanFunction}
\difficulty{3}
Let $m,n\in\mathbb N$, let $e$ be a bijection from the $m$ coordinates to the $n$
coordinates, let $f$ be a Boolean function on the sign cube of dimension $m$, and assign
to each of the $n$ coordinates one of three statuses: free, fixed to $1$, or fixed to
$-1$. The decision tree depth of the restriction of the reindexed function using these
statuses equals the decision tree depth of the restriction of $f$ using, at coordinate
$i$, the status assigned to $e(i)$.
\end{theorem}

\begin{theorem}[Reindexing the failure indicator]
\lean{FABL.coordSwitchingFailureIndicator_reindexBooleanFunction}
\label{FABL.coordSwitchingFailureIndicator_reindexBooleanFunction}
\leanok
\uses{FABL.BooleanFunction, FABL.CoordRestriction, FABL.coordSwitchingFailureIndicator, FABL.reindexBooleanFunction, FABL.restrictedDecisionTreeDepth_reindexBooleanFunction}
\difficulty{3}
Let $m,n$ be natural numbers, let $e:[m]\to[n]$ be a bijection, let
$f:\{-1,1\}^m\to\{\pm1\}$ be a Boolean function, let $k$ be a natural number, and
let $\rho:[n]\to\{\text{free},\text{fixed }1,\text{fixed }-1\}$ be a coordinate
restriction. Reindexing $f$ along $e$ and then restricting by $\rho$ gives the same
coordinate-switching failure indicator at threshold $k$ as restricting $f$ by the
restriction $i\mapsto\rho(e(i))$.
\end{theorem}

\begin{theorem}[Reindexing assignment weights]
\lean{FABL.restrictionAssignmentWeightAt_reindex}
\label{FABL.restrictionAssignmentWeightAt_reindex}
\leanok
\uses{FABL.CoordRestriction, FABL.coordRestrictionWeightAt, FABL.restrictionAssignmentWeightAt}
\difficulty{3}
Let $m,n$ be natural numbers, let $e$ be a bijection from the $m$ coordinates to the $n$
coordinates, let $b4$ be real, and assign a restriction state to each of the $n$
coordinates. Pulling the assignment back along $e$ leaves its product weight unchanged,
when each coordinate weight is $b4$ for a free state and $(1-b4)/2$ for either fixed
state.
\end{theorem}

\begin{theorem}[Reindexing preserves failure probability]
\lean{FABL.coordSwitchingFailureProbability_reindexBooleanFunction}
\label{FABL.coordSwitchingFailureProbability_reindexBooleanFunction}
\leanok
\uses{FABL.BooleanFunction, FABL.CoordRestriction, FABL.coordSwitchingFailureIndicator, FABL.coordSwitchingFailureIndicator_reindexBooleanFunction, FABL.coordSwitchingFailureProbability, FABL.reindexBooleanFunction, FABL.restrictionAssignmentWeightAt, FABL.restrictionAssignmentWeightAt_reindex}
\difficulty{4}
Let $m,n$ be natural numbers, let $e:[m]\to[n]$ be a bijection, let
$f:\{-1,1\}^m\to\{\pm1\}$ be a Boolean function, let $\delta\in\mathbb R$, and
let $k$ be a natural number. The coordinate-switching failure probability at parameter
$\delta$ and threshold $k$ is the same for $f$ and for the function on $[n]$ obtained
by reindexing its coordinates along $e$.
\end{theorem}

\begin{theorem}[Invariance under coordinate relabeling]
\lean{FABL.switchingFailureProbability_reindexBooleanFunction}
\label{FABL.switchingFailureProbability_reindexBooleanFunction}
\leanok
\uses{FABL.BooleanFunction, FABL.coordSwitchingFailureProbability, FABL.coordSwitchingFailureProbability_eq, FABL.coordSwitchingFailureProbability_reindexBooleanFunction, FABL.reindexBooleanFunction, FABL.switchingFailureProbability}
\difficulty{3}
Let $m,n\in\mathbb N$, let $e:[m]\to[n]$ be a bijection, let
$f:\{-1,1\}^m\to\{-1,1\}$ be a sign-valued Boolean function, and let
$\delta\in\mathbb R$ and $k\in\mathbb N$. Relabel the coordinates of $f$ using $e$
to obtain the function $g(x)=f(i\mapsto x(e(i)))$ on $\{-1,1\}^n$. Then
$P_{\delta,k}(g)=P_{\delta,k}(f)$.
\end{theorem}

\begin{definition}[Free coordinate embedding]
\lean{FABL.freeCoordinateEmbedding}
\statementsource{Analysis-of-Boolean-Functions-by-Ryan-ODonnell}{
  pdf-ffbe9c7b3306-p098-b003,
  pdf-ffbe9c7b3306-p098-b004
}
\label{FABL.freeCoordinateEmbedding}
\leanok
For a finite set $J\subseteq\{0,\ldots,n-1\}$, the free coordinate embedding sends the
$k$th element of $\{0,\ldots,|J|-1\}$ to the $k$th element of $J$ in increasing order.
\end{definition}

\begin{theorem}[Embedding of a finite coordinate set]
\lean{FABL.freeCoordinateEmbedding_equivFin}
\label{FABL.freeCoordinateEmbedding_equivFin}
\leanok
\uses{FABL.freeCoordinateEmbedding}
\difficulty{1}
Let $n$ be a nonnegative integer, let $J\subseteq\{0,\ldots,n-1\}$, and enumerate its
elements by the canonical bijection from $\{0,\ldots,|J|-1\}$ to $J$. The embedding of
these positions into $\{0,\ldots,n-1\}$ sends each position to the corresponding element
of $J$.
\end{theorem}

\begin{definition}[Reindexed sign restriction]
\lean{FABL.reindexedSignRestriction}
\statementsource{Analysis-of-Boolean-Functions-by-Ryan-ODonnell}{pdf-ffbe9c7b3306-p098-b004}
\label{FABL.reindexedSignRestriction}
\leanok
\uses{FABL.BooleanFunction, FABL.FixedSignCube, FABL.signRestriction}
Let $f:\{\pm1\}^n\to\{\pm1\}$, let $J\subseteq\{0,\ldots,n-1\}$, and assign signs
$z_i\in\{\pm1\}$ to every index outside $J$. The reindexed restriction is the function
on $\{\pm1\}^{|J|}$ obtained by inserting the input signs into the coordinates of $J$ in
increasing order, using $z_i$ on all coordinates outside $J$, and evaluating $f$.
\end{definition}

\begin{theorem}[Evaluation of a Reindexed Sign Restriction]
\lean{FABL.reindexedSignRestriction_apply}
\label{FABL.reindexedSignRestriction_apply}
\leanok
\uses{FABL.BooleanFunction, FABL.FixedSignCube, FABL.combineSignCube, FABL.reindexedSignRestriction}
\difficulty{1}
Let $n$ be a natural number, let $f\colon\{-1,1\}^n\to\{-1,1\}$, let $J$ be a set of
coordinates, and fix signs on the coordinates outside $J$. For every sign vector $x$ of
dimension $|J|$, the reindexed restriction of $f$ at $x$ equals $f$ evaluated on the
full sign vector formed from the corresponding coordinates of $x$ on $J$ and the fixed
signs outside $J$.
\end{theorem}

\begin{theorem}[Binary restriction via reindexed coordinates]
\lean{FABL.restrictedBinaryFunction_eq_comp_reindexedSignRestriction}
\label{FABL.restrictedBinaryFunction_eq_comp_reindexedSignRestriction}
\leanok
\uses{FABL.BooleanFunction, FABL.FixedSignCube, FABL.binaryCubeSignEquiv, FABL.freeCoordinateEmbedding, FABL.reindexedSignRestriction, FABL.restrictedBinaryFunction}
\difficulty{4}
Let $f$ be a function from the sign cube on $n$ coordinates to the signs, let $J$ be a
set of coordinates, and fix signs outside $J$. For every binary input on the full
coordinate set, evaluating the resulting restricted function is the same as converting
its coordinates in $J$ to signs, reindexing them by their positions in $J$, and
evaluating the corresponding reindexed restriction of $f$.
\end{theorem}

\begin{theorem}[Decision-tree depth under coordinate reindexing]
\lean{FABL.restrictedDecisionTreeDepth_eq_reindexedSignRestriction}
\label{FABL.restrictedDecisionTreeDepth_eq_reindexedSignRestriction}
\leanok
\uses{FABL.BooleanFunction, FABL.FixedSignCube, FABL.F₂DecisionTree.decisionTreeDepth, FABL.F₂DecisionTree.decisionTreeDepth_comp_embedding, FABL.binaryCubeSignEquiv, FABL.freeCoordinateEmbedding, FABL.reindexedSignRestriction, FABL.restrictedBinaryFunction_eq_comp_reindexedSignRestriction, FABL.restrictedDecisionTreeDepth}
\difficulty{3}
Let $f$ be a function from the sign cube on $n$ coordinates to the signs, let $J$ be a
set of coordinates, and fix signs outside $J$. The decision-tree depth of the binary
function obtained by restricting $f$ to those fixed signs equals the decision-tree depth
of the corresponding reindexed restriction, viewed as a function on the binary cube with
$|J|$ coordinates.
\end{theorem}

\begin{definition}[Composition of coordinate restrictions]
\lean{FABL.CoordRestriction.compose}
\label{FABL.CoordRestriction.compose}
\leanok
\uses{FABL.CoordRestriction}
For coordinate restrictions taking the values free, fixed at $1$, or fixed at $-1$,
define composition by letting a free first restriction pass through the second
restriction, while either fixed first restriction remains fixed at its chosen value
regardless of the second restriction.
\end{definition}

\begin{definition}[Pointwise composition of restrictions]
\lean{FABL.composeCoordRestrictions}
\label{FABL.composeCoordRestrictions}
\leanok
\uses{FABL.CoordRestriction, FABL.CoordRestriction.compose}
Let $n$ be a natural number, and let $r,s:[n]\to\{\text{free},\text{fixed at
}1,\text{fixed at }-1\}$ be coordinate restrictions. Their composition is the coordinate
restriction whose value at each coordinate is the composition of $r$ and $s$ there.
\end{definition}

\begin{definition}[Product weight of coordinate restrictions]
\lean{FABL.coordRestrictionProductWeight}
\label{FABL.coordRestrictionProductWeight}
\leanok
\uses{FABL.CoordRestriction, FABL.coordRestrictionWeightAt}
Let $\iota$ be a finite type, let $\delta\in\mathbb{R}$, and assign to each $i\in\iota$
a coordinate restriction $\rho(i)$. The product weight of this assignment is the
product, over all $i\in\iota$, of the weight at parameter $\delta$ of the restriction
$\rho(i)$.
\end{definition}

\begin{theorem}[Finite product weight on coordinates]
\lean{FABL.coordRestrictionProductWeight_fin}
\label{FABL.coordRestrictionProductWeight_fin}
\leanok
\uses{FABL.CoordRestriction, FABL.coordRestrictionProductWeight, FABL.restrictionAssignmentWeightAt}
\difficulty{1}
Let $n$ be a natural number, let $\delta\in\mathbb{R}$, and assign one of the states
free, fixed to $1$, or fixed to $-1$ to each of the $n$ coordinates. The product of the
coordinate weights over these coordinates equals the restriction assignment weight: each
free coordinate contributes $\delta$, and each fixed coordinate contributes
$(1-\delta)/2$.
\end{theorem}

\begin{theorem}[Total weight of coordinate restrictions]
\lean{FABL.sum_coordRestrictionProductWeight}
\label{FABL.sum_coordRestrictionProductWeight}
\leanok
\uses{FABL.CoordRestriction, FABL.coordRestrictionProductWeight, FABL.coordRestrictionWeightAt, FABL.sum_coordRestrictionWeightAt}
\difficulty{3}
Let $\iota$ be a finite set, and assign to each coordinate restriction one of the
weights $\delta$ for a free coordinate and $(1-\delta)/2$ for either fixed sign, for
any real $\delta$. The sum, over all assignments of restrictions to the coordinates in
$\iota$, of the product of their coordinate weights is $1$.
\end{theorem}

\begin{theorem}[Invariance of product weight under reindexing]
\lean{FABL.coordRestrictionProductWeight_reindex}
\label{FABL.coordRestrictionProductWeight_reindex}
\leanok
\uses{FABL.CoordRestriction, FABL.coordRestrictionProductWeight, FABL.coordRestrictionWeightAt}
\difficulty{3}
Let $I$ and $K$ be finite sets, let $e:I\to K$ be a bijection, let
$\delta\in\mathbb{R}$, and assign one of the states free, fixed to $1$, or fixed to
$-1$ to each element of $I$. The product of the corresponding coordinate weights is
unchanged when the assignment is transported to $K$ along $e$.
\end{theorem}

\begin{theorem}[Product weight across a coordinate split]
\lean{FABL.coordRestrictionProductWeight_split}
\label{FABL.coordRestrictionProductWeight_split}
\leanok
\uses{FABL.CoordRestriction, FABL.FixedIndex, FABL.coordRestrictionProductWeight, FABL.coordRestrictionWeightAt}
\difficulty{4}
Let $n$ be a natural number, let $J$ be a set of coordinates in $br^n$ indexed by
$	ext{Fin}(n)$, and assign a coordinate restriction to each index in $J$ and to each
index outside $J$. The product of the coordinate weights over all indices, using these
assignments on their respective parts, equals the product of the two separate weights
over $J$ and its complement.
\end{theorem}

\begin{theorem}[Restriction weights on selected coordinates]
\lean{FABL.sum_coordRestrictionProductWeight_restrict}
\label{FABL.sum_coordRestrictionProductWeight_restrict}
\leanok
\uses{FABL.CoordRestriction, FABL.FixedIndex, FABL.coordRestrictionProductWeight, FABL.coordRestrictionProductWeight_split, FABL.restrictionAssignmentWeightAt, FABL.sum_coordRestrictionProductWeight}
\difficulty{5}
Let $n$ be a nonnegative integer, let $\delta$ be real, let
$J\subseteq\{1,\ldots,n\}$, and let $g$ be a real-valued function on assignments of
coordinate restrictions to $J$. Sum over all assignments of restrictions to the $n$
coordinates, weighting each by the product of its coordinate weights and evaluating $g$
on the restriction to $J$. This sum equals the sum over assignments on $J$ alone of the
product of their coordinate weights times $g$ of that assignment.
\end{theorem}

\begin{theorem}[Restriction weight under coordinate projection]
\lean{FABL.sum_restrictionAssignmentWeightAt_comp_freeCoordinateEmbedding}
\label{FABL.sum_restrictionAssignmentWeightAt_comp_freeCoordinateEmbedding}
\leanok
\uses{FABL.CoordRestriction, FABL.coordRestrictionProductWeight, FABL.coordRestrictionProductWeight_fin, FABL.coordRestrictionProductWeight_reindex, FABL.freeCoordinateEmbedding, FABL.restrictionAssignmentWeightAt, FABL.sum_coordRestrictionProductWeight_restrict}
\difficulty{4}
Let $n$ be a nonnegative integer, let $J$ be a subset of the $n$ coordinates, let
$\delta\in\mathbb R$, and let $g$ be any real-valued function on restriction assignments
to $|J|$ coordinates. Under the product weights that assign weight $\delta$ to a free
coordinate and $(1-\delta)/2$ to either fixed value, the weighted sum over all
assignments on the full coordinate set of $g$ applied to their restrictions on $J$
equals the weighted sum of $g$ over all assignments on $J$.
\end{theorem}

\begin{definition}[Combining two restriction stages]
\lean{FABL.twoStageCoordRestriction}
\label{FABL.twoStageCoordRestriction}
\leanok
\uses{FABL.CoordRestriction, FABL.CoordRestriction.compose}
Let $n$ be a natural number and assign to each coordinate in $[n]$ an ordered pair of
coordinate restrictions, representing a first and a second stage. The resulting
restriction at each coordinate is the composition of its first-stage restriction with
its second-stage restriction.
\end{definition}

\begin{definition}[Weight of two restriction stages]
\lean{FABL.twoStageCoordRestrictionWeight}
\label{FABL.twoStageCoordRestrictionWeight}
\leanok
\uses{FABL.CoordRestriction, FABL.coordRestrictionWeightAt}
Let $n$ be a natural number, let $\delta_1,\delta_2\in\mathbb{R}$, and assign to each
coordinate in $[n]$ an ordered pair of coordinate restrictions. The two-stage weight is
the product over coordinates of the first-stage restriction's weight at parameter
$\delta_1$ multiplied by the second-stage restriction's weight at parameter $\delta_2$.
\end{definition}

\begin{theorem}[Composition of two coordinate restrictions]
\lean{FABL.sum_twoStageCoordRestrictionWeight_localFiber}
\label{FABL.sum_twoStageCoordRestrictionWeight_localFiber}
\leanok
\uses{FABL.CoordRestriction, FABL.CoordRestriction.compose, FABL.coordRestrictionWeightAt}
\difficulty{4}
For real parameters $\delta_1,\delta_2$ and a target restriction
$t\in\{\mathrm{free},\mathrm{fix\ }1,\mathrm{fix\ }(-1)\}$, sum over all ordered pairs
of coordinate restrictions whose composition is $t$, weighting a pair $(r_1,r_2)$ by the
product of its weights at $\delta_1$ and $\delta_2$. This sum equals the weight of $t$
at $\delta_1\delta_2$.
\end{theorem}

\begin{theorem}[Two-stage restriction weight over a fiber]
\lean{FABL.sum_twoStageCoordRestrictionWeight_fiber}
\label{FABL.sum_twoStageCoordRestrictionWeight_fiber}
\leanok
\uses{FABL.CoordRestriction, FABL.CoordRestriction.compose, FABL.coordRestrictionWeightAt, FABL.restrictionAssignmentWeightAt, FABL.sum_twoStageCoordRestrictionWeight_localFiber, FABL.twoStageCoordRestriction, FABL.twoStageCoordRestrictionWeight}
\difficulty{5}
Let $n$ be a natural number, let $\delta_1,\delta_2\in\mathbb R$, and fix a
coordinate restriction assignment on $n$ coordinates. Sum the two-stage weights over all
pairs of restriction assignments whose coordinatewise composition is the fixed
assignment. This sum equals the weight of the fixed assignment with parameter
$\delta_1\delta_2$.
\end{theorem}

\begin{theorem}[Composition of weighted restrictions]
\lean{FABL.sum_restrictionAssignmentWeightAt_compose}
\label{FABL.sum_restrictionAssignmentWeightAt_compose}
\leanok
\uses{FABL.CoordRestriction, FABL.composeCoordRestrictions, FABL.restrictionAssignmentWeightAt, FABL.sum_twoStageCoordRestrictionWeight_fiber, FABL.twoStageCoordRestriction, FABL.twoStageCoordRestrictionWeight}
\difficulty{5}
Let $n$ be a nonnegative integer, let $\delta_1,\delta_2\in\mathbb R$, and let $g$ be
any real-valued function on assignments of coordinate restrictions to the $n$
coordinates. Compose two assignments coordinatewise, with a first-stage fixed value
taking precedence and a free first-stage coordinate inheriting the second-stage
restriction. The weighted sum over the two stages, using parameters $\delta_1$ and
$\delta_2$, of $g$ evaluated on the composed assignment equals the weighted sum of $g$
over single-stage assignments with parameter $\delta_1\delta_2$.
\end{theorem}

\begin{theorem}[Composition of coordinate restrictions]
\lean{FABL.complete_compose_coordRestrictionOf}
\label{FABL.complete_compose_coordRestrictionOf}
\leanok
\uses{FABL.CoordRestriction, FABL.CoordRestriction.complete, FABL.CoordRestriction.compose, FABL.FixedSignCube, FABL.combineSignCube, FABL.combineSignCube_apply_fixed, FABL.combineSignCube_apply_free, FABL.composeCoordRestrictions, FABL.coordRestrictionOf, FABL.freeCoordinateEmbedding}
\difficulty{4}
Let $n$ be a natural number, let $J$ be a subset of the coordinates $	ext{Fin}(n)$, let
$z$ assign a sign in ${ 1,1}$ to each coordinate outside $J$, let $
ho$ be a
coordinate restriction on $	ext{Fin}(n)$, and let $x	ext{in }F_2^n$. Completing the
coordinatewise composition of the restriction determined by $J,z$ with $
ho$ at $x$
gives the sign vector whose coordinates in $J$ are obtained by completing the
restriction reindexed to $J$ at the corresponding reindexed coordinates of $x$, and
whose coordinates outside $J$ are given by $z$.
\end{theorem}

\begin{theorem}[Composition of coordinate restrictions]
\lean{FABL.restrict_compose_coordRestrictionOf}
\label{FABL.restrict_compose_coordRestrictionOf}
\leanok
\uses{FABL.BooleanFunction, FABL.CoordRestriction, FABL.CoordRestriction.restrict, FABL.FixedSignCube, FABL.complete_compose_coordRestrictionOf, FABL.composeCoordRestrictions, FABL.coordRestrictionOf, FABL.freeCoordinateEmbedding, FABL.reindexedSignRestriction, FABL.reindexedSignRestriction_apply}
\difficulty{3}
Let $f$ be a function from the sign cube on $n$ coordinates to the signs, let $J$ be a
set of coordinates, and fix signs $z$ on the complement of $J$. Let $\rho$ assign to
each coordinate one of three rules: leave it free, fix it to $1$, or fix it to $-1$.
Applying the fixed assignment $z$ outside $J$ and then the rules $\rho$ agrees with
first restricting $f$ to the coordinates in $J$ using $z$, reindexing those free
coordinates by $J$, and applying at each reindexed coordinate the rule assigned by
$\rho$ to its corresponding original coordinate.
\end{theorem}

\begin{theorem}[Depth under composed coordinate restrictions]
\lean{FABL.decisionTreeDepth_restrict_compose_coordRestrictionOf}
\label{FABL.decisionTreeDepth_restrict_compose_coordRestrictionOf}
\leanok
\uses{FABL.BooleanFunction, FABL.CoordRestriction, FABL.CoordRestriction.restrict, FABL.FixedSignCube, FABL.F₂DecisionTree.decisionTreeDepth, FABL.F₂DecisionTree.decisionTreeDepth_comp_embedding, FABL.composeCoordRestrictions, FABL.coordRestrictionOf, FABL.freeCoordinateEmbedding, FABL.reindexedSignRestriction, FABL.restrict_compose_coordRestrictionOf}
\difficulty{3}
Let $n\in\mathbb{N}$, let $f:\{−1,1\}^n\to\{−1,1\}$, let $J\subseteq\{1,\ldots,n\}$, and
fix signs on all coordinates outside $J$. For each coordinate, choose a second
restriction state: free, fixed to $1$, or fixed to $-1$. Compose this state assignment
with the assignment that leaves coordinates in $J$ free and fixes the other coordinates
to the chosen signs; the resulting restricted function on the binary cube has the same
minimum decision tree depth as the function obtained by first restricting $f$ to the
fixed signs outside $J$, reindexing the free coordinates by $J$, and then applying the
second restriction on those coordinates.
\end{theorem}

\begin{theorem}[Switching failure under composed restrictions]
\lean{FABL.coordSwitchingFailureIndicator_compose_coordRestrictionOf}
\label{FABL.coordSwitchingFailureIndicator_compose_coordRestrictionOf}
\leanok
\uses{FABL.BooleanFunction, FABL.CoordRestriction, FABL.FixedSignCube, FABL.composeCoordRestrictions, FABL.coordRestrictionOf, FABL.coordSwitchingFailureIndicator, FABL.decisionTreeDepth_restrict_compose_coordRestrictionOf, FABL.freeCoordinateEmbedding, FABL.reindexedSignRestriction}
\difficulty{3}
Let $f:\{\pm1\}^n\to\{\pm1\}$ be a Boolean function, let $k$ be a natural number, let
$J\subseteq\text{Fin}(n)$, let $z$ assign a sign to every coordinate outside $J$, and
let $\sigma$ assign a coordinate restriction to each coordinate. First impose the fixed
signs $z$ outside $J$, leaving the coordinates in $J$ free, and then compose this
assignment with $\sigma$. The switching failure indicator of $f$ at threshold $k$ for
the resulting assignment equals the indicator for the function on the free coordinates
obtained by fixing the outside coordinates to $z$, at threshold $k$, with restrictions
given by $\sigma$ on the coordinates in $J$ via their increasing enumeration.
\end{theorem}

\begin{theorem}[Restriction identity via coordinate assignments]
\lean{FABL.switchingFailureProbability_reindexedSignRestriction}
\label{FABL.switchingFailureProbability_reindexedSignRestriction}
\leanok
\uses{FABL.BooleanFunction, FABL.CoordRestriction, FABL.FixedSignCube, FABL.composeCoordRestrictions, FABL.coordRestrictionOf, FABL.coordSwitchingFailureIndicator, FABL.coordSwitchingFailureIndicator_compose_coordRestrictionOf, FABL.coordSwitchingFailureProbability, FABL.coordSwitchingFailureProbability_eq, FABL.reindexedSignRestriction, FABL.restrictionAssignmentWeightAt, FABL.sum_restrictionAssignmentWeightAt_comp_freeCoordinateEmbedding, FABL.switchingFailureProbability}
\difficulty{4}
Let $f:\{-1,1\}^n\to\{-1,1\}$ be a sign-valued Boolean function, let
$J\subseteq[n]$, let $z\in\{-1,1\}^{[n]\setminus J}$ specify the fixed signs
outside $J$, and let $\delta\in\mathbb R$ and $k\in\mathbb N$. The
switching-failure probability at parameter $\delta$ and threshold $k$ for the
restriction of $f$ to $J$ equals the sum over all coordinate assignments
$\rho:[n]\to\{\text{free},+1,-1\}$ of their product weights times the indicator
that the function obtained by composing the restriction specified by $(J,z)$ with
$\rho$ has decision-tree depth at least $k$. Each coordinate's weight is $\delta$ when
free and $(1-\delta)/2$ when fixed to either sign.
\end{theorem}

\begin{definition}[Restriction atom from an assignment]
\lean{FABL.restrictionAtomOfAssignment}
\statementsource{lmn-constant-depth-fourier}{pdf-63759bd909d5-p005-b001}
\statementsource{Analysis-of-Boolean-Functions-by-Ryan-ODonnell}{
  pdf-ffbe9c7b3306-p098-b003,
  pdf-ffbe9c7b3306-p098-b004
}
\label{FABL.restrictionAtomOfAssignment}
\leanok
\uses{FABL.CoordRestriction, FABL.RandomRestrictionAtom, FABL.randomRestrictionAtomEquiv}
Given a coordinate assignment that marks each of the $n$ coordinates as free, fixed to
$+1$, or fixed to $-1$, its restriction atom is the pair consisting of the set $J$ of
coordinates marked free and the assignment of fixed signs to the coordinates outside
$J$.
\end{definition}

\begin{theorem}[Restriction recovered from an assignment]
\lean{FABL.coordRestrictionOf_restrictionAtomOfAssignment}
\label{FABL.coordRestrictionOf_restrictionAtomOfAssignment}
\leanok
\uses{FABL.CoordRestriction, FABL.coordRestrictionOf, FABL.randomRestrictionAtomEquiv, FABL.restrictionAtomOfAssignment}
\difficulty{1}
Let $n$ be a natural number and let $\rho$ assign to each of the $n$ coordinates one of
three states: free, fixed to $1$, or fixed to $-1$. The restriction atom associated to
$\rho$ determines exactly the coordinate restriction $\rho$.
\end{theorem}

\begin{theorem}[Random restriction expectation as an assignment sum]
\lean{FABL.expectRandomRestriction_eq_sum_assignment}
\label{FABL.expectRandomRestriction_eq_sum_assignment}
\leanok
\uses{FABL.CoordRestriction, FABL.FixedSignCube, FABL.RandomRestrictionAtom, FABL.deltaRandomSubsetWeight, FABL.expectRandomRestriction, FABL.randomRestrictionAtomEquiv, FABL.restrictionAssignmentWeightAt, FABL.restrictionAssignmentWeightAt_coordRestrictionOf_eq_div, FABL.restrictionAtomOfAssignment}
\difficulty{5}
Let $n$ be a nonnegative integer, let $\delta\in\mathbb R$, and let $g$ be any
real-valued function of a free-coordinate subset of $\{1,\ldots,n\}$ and signs on its
complement. The random restriction expectation of $g$ equals the sum over all coordinate
assignments in which each coordinate is free, fixed to $+1$, or fixed to $-1$: each
assignment contributes its value of $g$ multiplied by the product of coordinate weights,
where a free coordinate has weight $\delta$ and either fixed sign has weight
$(1-\delta)/2$.
\end{theorem}

\begin{theorem}[Algebraic restriction composition identity]
\lean{FABL.switchingFailureProbability_mul_algebraic}
\label{FABL.switchingFailureProbability_mul_algebraic}
\leanok
\uses{FABL.BooleanFunction, FABL.coordRestrictionOf_restrictionAtomOfAssignment, FABL.coordSwitchingFailureIndicator, FABL.coordSwitchingFailureProbability, FABL.coordSwitchingFailureProbability_eq, FABL.expectRandomRestriction, FABL.expectRandomRestriction_eq_sum_assignment, FABL.reindexedSignRestriction, FABL.sum_restrictionAssignmentWeightAt_compose, FABL.switchingFailureProbability, FABL.switchingFailureProbability_reindexedSignRestriction}
\difficulty{4}
Let $f:\{-1,1\}^n\to\{-1,1\}$ be a sign-valued Boolean function, let
$\delta_1,\delta_2\in\mathbb R$, and let $k\in\mathbb N$. The switching-failure
probability of $f$ at parameter $\delta_1\delta_2$ and threshold $k$ equals the
average over first restrictions at parameter $\delta_1$ of the switching-failure
probabilities at parameter $\delta_2$ and threshold $k$ for the corresponding functions
on the free coordinates.
\end{theorem}

\begin{theorem}[Multiplicativity under random restrictions]
\lean{FABL.switchingFailureProbability_mul}
\label{FABL.switchingFailureProbability_mul}
\leanok
\uses{FABL.BooleanFunction, FABL.expectRandomRestriction, FABL.reindexedSignRestriction, FABL.switchingFailureProbability, FABL.switchingFailureProbability_mul_algebraic}
\difficulty{3}
Let $f:\{-1,1\}^n\to\{-1,1\}$ be a sign-valued Boolean function, let $k\in\mathbb
N$, and let $0\le\delta_1,\delta_2\le 1$. Then $P_{\delta_1\delta_2,k}(f)$ equals
the average, over a first random restriction at parameter $\delta_1$, of the
switching-failure probability at parameter $\delta_2$ and threshold $k$ for the
resulting function on its free coordinates.
\end{theorem}
